\documentclass{article}
\usepackage[margin=0.8in]{geometry}
\usepackage{amsmath,amssymb,amsthm,mathtools}
\usepackage{mathrsfs,bm,bbm}
\usepackage{graphicx,booktabs,array,multirow,tabularx,longtable}
\usepackage{enumitem}
\usepackage{xcolor}
\usepackage{algorithm,algpseudocode}
\usepackage{subcaption,tikz}
\usepackage{siunitx,cancel,accents}
\usepackage{microtype}
\usepackage[numbers,sort&compress]{natbib}
\usepackage[colorlinks=true,linkcolor=blue,citecolor=blue,urlcolor=blue]{hyperref}
\usepackage{cleveref}
\setlength{\emergencystretch}{2em}
\newtheorem{assumption}{Assumption}
\newtheorem{definition}{Definition}
\newtheorem{theorem}{Theorem}
\newtheorem{lemma}{Lemma}
\newtheorem{proposition}{Proposition}
\newtheorem{openendedquestion}{Open-ended Question}
\newtheorem{conjecture}{Conjecture}
\newcommand{\coreref}[1]{\ref{#1}}

\begin{document}

\sloppy

\section*{eid\_\allowbreak{}hidden\_\allowbreak{}ordinal\_\allowbreak{}eif\_\allowbreak{}operator}

\noindent\textit{Specialization:} v1\par

\paragraph{TL;DR.} For fixed \(K\geq3\), the exact continuous-outcome canonical gradient remains the product-bridge score minus the pointwise \(K^2\)-cell projection and one conditional Gram correction.  The implementable estimator now fits one raw product law, threshold-cleans and renormalizes both measurement matrices without a support oracle or beta-min condition, and forms every bridge and operator from that cleaned law.  If the raw uniform errors are \(o_P(\tau_{S,n})\), true-zero cells are recovered exactly; under the displayed cleaned product rates, cross-fitting attains the efficient joint Gaussian and maximum-\(t\) limits.  Proving the named global spline product-likelihood rates remains an explicit open question.

\section{Project justification}

\paragraph{Gap.} The binary hidden-treatment literature characterizes efficiency but computes its continuous-outcome correction through a partition-dependent sequence, while categorical product-bridge work establishes validity without observed-data canonicality. Existing continuous three-proxy identification results also do not provide the measurable uniformly bounded differential inverse required to identify the full observed DQM tangent closure.

\paragraph{Niche.} Fixed latent cardinality, quantitative three-view separation, common outcome support, and paired categorical measurements permit both exact inversion of the differentiated factor map and elimination of the infinite-dimensional outcome nuisance before a finite conditional Gram solve.

\paragraph{Fill.} The differential-score lifting theorem closes tangent completeness and yields the exact finite canonical projector.  For implementation, one label-aligned raw product fit is thresholded and column-renormalized before any bridge, posterior, Gram, or score is constructed.  The cleaning algebra proves exact true-zero recovery from \(r_{A,\infty}+r_{Z,\infty}=o_P(\tau_{S,n})\), with cleaned errors \(r_A^c=r_A+\tau_{S,n}\) and \(r_Z^c=r_Z+\tau_{S,n}\), so the coherent mean bias is quadratic and contains no additive leakage allowance.  The Gram inverse still uses the computable cutoff \(\varepsilon/2\).

\section{Related work}

Zhou and Tchetgen Tchetgen (2026) are the closest theorem-level comparator: their Theorems 2--4 treat the binary observed tangent space and influence functions, and Supplement S8.2 obtains the continuous-outcome efficiency bound through a refining outcome partition.  The present quantitatively separated multilevel specialization instead uses the exact \(D\)-orthogonal projector \(I-D^{-1}B(B^{\mathsf T}D^{-1}B)^+B^{\mathsf T}\) and one conditional Gram solve.  Bian, Zhou, and Cui (2026) supply the categorical inverse-feature product-bridge precursor, but not observed-data canonicality or the support-cleaned projected-score remainder.  Allman, Matias, and Rhodes (2009), Hu and Schennach (2008), and Deaner (2023) inform latent-factor identification; none supplies the measurable bounded differential inverse proved here.  Bonhomme, Jochmans, and Robin (2016) and Tsybakov (2009) motivate coherent repeated-measurement estimation and the smooth sieve benchmark, but they do not prove the open curvature and uniform-rate properties for the specific product likelihood used here.

\section{Setup}

Target (estimand / structural parameter): \(\theta=C\psi=C\bigl(E\{Y(1)\},\ldots,E\{Y(K)\}\bigr)^{\mathsf T}\).
Primary functional / estimator / diagnostic: \(\theta=C\,E_P\bigl[\bigl\{\int_{\mathcal Y}y f_j(y\mid X)d\nu(y)\bigr\}_{j=1}^{K}\bigr]\).
\begin{itemize}
  \item \texttt{K} --- \(\mathbb N\); \texttt{Fixed number of latent treatment levels and categories in each measurement view.}
  \item \texttt{calJ} --- \texttt{finite set}; \(\{1,\ldots,K\}\); \texttt{Latent ordered treatment state space.}
  \par\smallskip\noindent\emph{Definition.} \(\mathcal J=\{1,\ldots,K\}\)
  \item \texttt{calA} --- \texttt{finite set}; \(\{1,\ldots,K\}\); \texttt{Surrogate state space.}
  \par\smallskip\noindent\emph{Definition.} \(\mathcal A=\{1,\ldots,K\}\)
  \item \texttt{calZ} --- \texttt{finite set}; \(\{1,\ldots,K\}\); \texttt{Proxy state space.}
  \par\smallskip\noindent\emph{Definition.} \(\mathcal Z=\{1,\ldots,K\}\)
  \item \texttt{calX} --- \texttt{standard Borel space}; \texttt{Covariate space.}
  \item \texttt{calY} --- \texttt{measurable subset}; \(\mathbb R\); \texttt{Common outcome support.}
  \item \texttt{nu} --- \texttt{sigma-\allowbreak{}finite measure}; \(\mathcal Y\); \texttt{Dominating measure for class outcome laws.}
  \item \texttt{L} --- \texttt{positive real}; \((0,\infty)\); \texttt{Outcome envelope.}
  \item \texttt{epsilon} --- \texttt{positive real}; \((0,1)\); \texttt{Quantitative regular-\allowbreak{}stratum margin.}
  \item \texttt{alpha} --- \texttt{probability level}; \((0,1)\); \texttt{Simultaneous inference error probability.}
  \item \texttt{q} --- \texttt{positive integer}; \texttt{Number of target contrasts.}
  \item \texttt{n} --- \texttt{positive integer}; \texttt{Sample size.}
  \item \texttt{Ffold} --- \texttt{positive integer}; \texttt{Fixed number of cross-\allowbreak{}fitting folds.}
  \item \texttt{fold} --- \texttt{deterministic fold map}; \(\{1,\ldots,n\}\to\{1,\ldots,F_{\mathrm{fold}}\}\); \texttt{Cross-\allowbreak{}fitting index map.}
  \par\smallskip\noindent\emph{Definition.} \(\mathrm{fold}(i)\) is the held-out fold containing unit \(i\)
  \item \texttt{J} --- \texttt{random variable}; \(\mathcal J\); \texttt{Hidden ordered treatment.}
  \item \texttt{X} --- \texttt{random variable}; \(\mathcal X\); \texttt{Observed pretreatment covariates.}
  \item \texttt{Y} --- \texttt{random variable}; \(\mathcal Y\); \texttt{Observed bounded continuous outcome.}
  \item \texttt{A} --- \texttt{random variable}; \(\mathcal A\); Observed categorical surrogate for \(J\).
  \item \texttt{Z} --- \texttt{random variable}; \(\mathcal Z\); Observed categorical auxiliary proxy for \(J\).
  \item \texttt{O} --- \texttt{random element}; \(\mathcal Y\times\mathcal A\times\mathcal Z\times\mathcal X\); \texttt{Observed unit.}
  \par\smallskip\noindent\emph{Definition.} \(O=(Y,A,Z,X)\)
  \item \texttt{P} --- \texttt{probability law}; laws of \(O\); \texttt{Observed-\allowbreak{}data law induced by a latent product factorization.}
  \item \texttt{PX} --- \texttt{probability law}; laws on \(\mathcal X\); \texttt{Covariate marginal law.}
  \par\smallskip\noindent\emph{Definition.} \(P_X\) is the \(X\)-marginal of \(P\)
  \item \texttt{Mclass} --- \texttt{model class}; \texttt{Evaluation model for identification and efficiency.}
  \par\smallskip\noindent\emph{Definition.} \(\mathcal M_{K,\varepsilon}\) is the regular hidden ordered-treatment class of laws \(P\)
  \item \texttt{Ypot} --- \texttt{potential-\allowbreak{}outcome vector}; \(\mathcal Y^{K}\); Potential outcomes under interventions on \(J\).
  \par\smallskip\noindent\emph{Definition.} \(Y^{\mathrm{pot}}=(Y(1),\ldots,Y(K))\)
  \item \texttt{pi} --- \texttt{measurable simplex-\allowbreak{}valued function}; \(\mathcal X\to\Delta^{K-1}\); \texttt{Latent treatment propensity.}
  \par\smallskip\noindent\emph{Definition.} \(\pi_j(x)=P(J=j\mid X=x)\)
  \item \texttt{f} --- \texttt{vector of conditional densities}; \(\mathcal Y\times\mathcal X\to\mathbb R_{+}^{K}\); \texttt{Class-\allowbreak{}specific outcome densities.}
  \par\smallskip\noindent\emph{Definition.} \(f_j(y\mid x)=dP(Y\in dy\mid J=j,X=x)/d\nu(y)\)
  \item \texttt{MA} --- \texttt{measurable matrix-\allowbreak{}valued function}; \(\mathcal X\to[0,1]^{K\times K}\); \texttt{Surrogate measurement factor.}
  \par\smallskip\noindent\emph{Definition.} \(M_A(x)_{aj}=P(A=a\mid J=j,X=x)\)
  \item \texttt{MZ} --- \texttt{measurable matrix-\allowbreak{}valued function}; \(\mathcal X\to[0,1]^{K\times K}\); \texttt{Proxy measurement factor.}
  \par\smallskip\noindent\emph{Definition.} \(M_Z(x)_{zj}=P(Z=z\mid J=j,X=x)\)
  \item \texttt{h0} --- \texttt{known bounded measurable feature map}; \(\mathcal Y\to\mathbb R^{K}\); \texttt{Outcome feature used only for factor recovery.}
  \item \texttt{Creg} --- \texttt{positive constant}; \texttt{Uniform regular-\allowbreak{}stratum bound.}
  \par\smallskip\noindent\emph{Definition.} \(C_{\mathrm{reg}}\) depends only on \(K\), \(L\), \(\varepsilon\), and the fixed feature envelope \(\|h_0\|_{\infty}\)
  \item \texttt{GammaY} --- \texttt{measurable matrix-\allowbreak{}valued function}; \(\mathcal X\to\mathbb R^{K\times K}\); \texttt{Outcome-\allowbreak{}feature factor.}
  \par\smallskip\noindent\emph{Definition.} \(\Gamma_Y(x)_{\cdot j}=E\{h_0(Y)\mid J=j,X=x\}\)
  \item \texttt{Pstar} --- \texttt{probability law}; laws of \(O\); \texttt{Explicit positive three-\allowbreak{}class continuous witness.}
  \item \texttt{m} --- \texttt{measurable vector-\allowbreak{}valued function}; \(\mathcal X\to\mathbb R^{K}\); \texttt{Class-\allowbreak{}specific outcome regression.}
  \par\smallskip\noindent\emph{Definition.} \(m_j(x)=\int_{\mathcal Y}y f_j(y\mid x)d\nu(y)\)
  \item \texttt{psi} --- \texttt{parameter vector}; \(\mathbb R^{K}\); \texttt{Mean potential outcomes identified after label recovery.}
  \par\smallskip\noindent\emph{Definition.} \(\psi_j=E_P\{m_j(X)\}\)
  \item \texttt{C} --- \texttt{fixed matrix}; \(\mathbb R^{q\times K}\); \texttt{Declared full-\allowbreak{}row-\allowbreak{}rank contrast matrix.}
  \item \texttt{theta} --- \texttt{parameter vector}; \(\mathbb R^{q}\); \texttt{Target ordered-\allowbreak{}treatment contrast.}
  \par\smallskip\noindent\emph{Definition.} \(\theta=C\psi\)
  \item \texttt{g} --- \texttt{conditional density}; \(\mathcal Y\times\mathcal X\to\mathbb R_{+}\); \texttt{Conditional mixture outcome density.}
  \par\smallskip\noindent\emph{Definition.} \(g(y\mid x)=\sum_{j=1}^{K}\pi_j(x)f_j(y\mid x)\)
  \item \texttt{w} --- \texttt{posterior weight vector}; \(\mathcal Y\times\mathcal X\to\Delta^{K-1}\); Posterior class weight given \((Y,X)\).
  \par\smallskip\noindent\emph{Definition.} \(w_j(y,x)=\pi_j(x)f_j(y\mid x)/g(y\mid x)\)
  \item \texttt{W} --- \texttt{diagonal matrix-\allowbreak{}valued function}; \(\mathcal Y\times\mathcal X\to\mathbb R^{K\times K}\); \texttt{Diagonal posterior-\allowbreak{}class weight matrix used in the exact Bayes reparameterization.}
  \par\smallskip\noindent\emph{Definition.} \(W(y,x)=\operatorname{diag}\{w(y,x)\}\)
  \item \texttt{B} --- \texttt{measurable matrix-\allowbreak{}valued function}; \(\mathcal X\to\mathbb R^{K^{2}\times K}\); \texttt{Khatri-\allowbreak{}-\allowbreak{}Rao measurement factor.}
  \par\smallskip\noindent\emph{Definition.} \(B_{\cdot j}(x)=M_A(x)_{\cdot j}\otimes M_Z(x)_{\cdot j}\)
  \item \texttt{d} --- \texttt{conditional cell-\allowbreak{}probability vector}; \(\mathcal Y\times\mathcal X\to\Delta^{K^{2}-1}\); Observed \((A,Z)\) cell probabilities given \((Y,X)\).
  \par\smallskip\noindent\emph{Definition.} \(d_{az}(y,x)=\sum_{j=1}^{K}w_j(y,x)B_{(a,z),j}(x)\)
  \item \texttt{D} --- \texttt{diagonal matrix-\allowbreak{}valued function}; \(\mathcal Y\times\mathcal X\to\mathbb R^{K^{2}\times K^{2}}\); \texttt{Cell-\allowbreak{}weight matrix.}
  \par\smallskip\noindent\emph{Definition.} \(D(y,x)=\operatorname{diag}\{d(y,x)\}\)
  \item \texttt{R} --- \texttt{posterior matrix-\allowbreak{}valued function}; \(\mathcal Y\times\mathcal X\to[0,1]^{K^{2}\times K}\); \texttt{Full observed posterior by cell,\allowbreak{} with its exact Bayes matrix factorization.}
  \par\smallskip\noindent\emph{Definition.} \(R_{(a,z),j}(y,x)=P(J=j\mid Y=y,A=a,Z=z,X=x)\) and, cellwise, \(R(y,x)=D(y,x)^{-1}B(x)W(y,x)\)
  \item \texttt{e} --- \texttt{canonical basis map}; \(r\mapsto e_r\in\mathbb R^K\); \texttt{One-\allowbreak{}hot basis vectors for latent and measurement categories.}
  \item \texttt{bridgeA} --- \texttt{vector of observed bridge functions}; \(\mathcal A\times\mathcal X\to\mathbb R^{K}\); \texttt{Surrogate inverse-\allowbreak{}feature bridge.}
  \par\smallskip\noindent\emph{Definition.} \(a_j(A,X)=e_j^{\mathsf T}M_A(X)^{-1}e_A\)
  \item \texttt{bridgeZ} --- \texttt{vector of observed bridge functions}; \(\mathcal Z\times\mathcal X\to\mathbb R^{K}\); \texttt{Proxy inverse-\allowbreak{}feature bridge.}
  \par\smallskip\noindent\emph{Definition.} \(z_j(Z,X)=e_j^{\mathsf T}M_Z(X)^{-1}e_Z\)
  \item \texttt{U} --- \texttt{uncentered score vector}; \(L^{2}(P)^{K}\); \texttt{Raw product-\allowbreak{}bridge score.}
  \par\smallskip\noindent\emph{Definition.} \(U_j=m_j(X)+a_j(A,X)z_j(Z,X)\{Y-m_j(X)\}/\pi_j(X)\)
  \item \texttt{phi} --- \texttt{influence-\allowbreak{}function vector}; \(L^{2}_{0}(P)^{K}\); \texttt{Raw product-\allowbreak{}bridge influence function.}
  \par\smallskip\noindent\emph{Definition.} \(\phi=U-\psi\)
  \item \texttt{Cadj} --- \texttt{finite contrast class}; \(\mathbb R^{K}\); \texttt{Adjacent maternal-\allowbreak{}smoking intensity step contrasts.}
  \par\smallskip\noindent\emph{Definition.} \(\mathcal C_{\mathrm{adj}}=\{c_j=e_{j+1}-e_j:j=1,\ldots,K-1\}\)
  \item \texttt{vAdj} --- \texttt{family of observable cell-\allowbreak{}score vectors}; \(\mathcal Y\times\mathcal X\to\mathbb R^{K^{2}}\); \texttt{Recovered cell representation of each adjacent-\allowbreak{}intensity raw contrast score.}
  \par\smallskip\noindent\emph{Definition.} For \(c_j\in\mathcal C_{\mathrm{adj}}\), \(v_j(y,x)=\{c_j^{\mathsf T}\phi(y,a,z,x)\}_{a,z}\)
  \item \texttt{Slat} --- \texttt{linear score class}; \(L^{2}_{0}(P)\); \texttt{Observed image of additive latent factor scores.}
  \par\smallskip\noindent\emph{Definition.} \(\mathcal S_{\mathrm{lat}}=\{E[s_X(X)+s_\pi(J,X)+s_Y(Y,J,X)+s_A(A,J,X)+s_Z(Z,J,X)\mid O]\}\), with every factor score centered under its own conditional factor law
  \item \texttt{T} --- \texttt{closed linear subspace}; \(L^{2}_{0}(P)\); \texttt{Full observed-\allowbreak{}data tangent space.}
  \par\smallskip\noindent\emph{Definition.} \(\mathcal T\) is the closure of scores of all observed-data DQM paths through \(P\) in the model
  \item \texttt{N} --- \texttt{pointwise linear operator}; \(\mathbb R^{K^{2}}\to\mathbb R^{K^{2}}\); The correct \(D\)-orthogonal cell projector onto \(\ker(B^{\mathsf T})\), written both in measurement-factor and posterior coordinates.
  \par\smallskip\noindent\emph{Definition.} \(N=N_B:=I_{K^{2}}-D^{-1}B(B^{\mathsf T}D^{-1}B)^{+}B^{\mathsf T}=I_{K^{2}}-R(R^{\mathsf T}DR)^{+}R^{\mathsf T}D\), and \((Nf_0)(O)=[N(Y,X)\{f_0(Y,a,z,X)\}_{a,z}]_{(A,Z)}\)
  \item \texttt{K3audit} --- \texttt{three-\allowbreak{}class algebraic test configuration}; Exact symbolic nonconstant-posterior audit of the first cell elimination at \(K=3\).
  \par\smallskip\noindent\emph{Definition.} \(\mathfrak A_3=(w,W,B,D,R,N_B)\) with \(w\in(0,1)^3\), \(\mathbf 1_3^{\mathsf T}w=1\), the \(w_j\) not all equal, \(W=\operatorname{diag}(w_1,w_2,w_3)\), \(B\in\mathbb R^{9\times3}\) of full column rank, \(Bw>0\) entrywise, \(D=\operatorname{diag}(Bw)\), \(R=D^{-1}BW\), and \(N_B=I_9-D^{-1}B(B^{\mathsf T}D^{-1}B)^{-1}B^{\mathsf T}\)
  \item \texttt{cellc} --- \texttt{vector of cell functions}; \(L^{2}(P)^{2K^{2}}\); \texttt{Cell representations of measurement-\allowbreak{}factor constraints.}
  \par\smallskip\noindent\emph{Definition.} \(c_{(j,a)}=R_{\cdot j}\mathbf 1\{A=a\}\) and \(c_{(j,z)}=R_{\cdot j}\mathbf 1\{Z=z\}\)
  \item \texttt{h} --- \texttt{residualized cell-\allowbreak{}function vector}; \(L^{2}(P)^{2K^{2}}\); Measurement-factor constraint vectors residualized by the correct \(D\)-orthogonal outcome-elimination projector.
  \par\smallskip\noindent\emph{Definition.} \(h_r(O)=[N(Y,X)c_r(Y,X)]_{(A,Z)}\) for each of the \(2K^{2}\) coordinates \(c_r\)
  \item \texttt{G} --- \texttt{conditional Gram matrix}; \(\mathcal X\to\mathbb R^{2K^{2}\times2K^{2}}\); \texttt{Conditional Gram operator of the correctly residualized measurement directions.}
  \par\smallskip\noindent\emph{Definition.} \(G(X)=E_P(hh^{\mathsf T}\mid X)\)
  \item \texttt{bf} --- \texttt{conditional coefficient map}; \(f_0\mapsto b_{f_0}\); \texttt{Right-\allowbreak{}hand side of the conditional Gram solve.}
  \par\smallskip\noindent\emph{Definition.} \(b_{f_0}(X)=E_P\{h(O)(Nf_0)(O)\mid X\}\)
  \item \texttt{bU} --- \texttt{conditional coefficient matrix}; \(\mathcal X\to\mathbb R^{2K^{2}\times K}\); \texttt{Gram right-\allowbreak{}hand side for the product-\allowbreak{}bridge score.}
  \par\smallskip\noindent\emph{Definition.} \(b_U(X)=E_P\{h(O)(NU)(O)^{\mathsf T}\mid X\}\)
  \item \texttt{LambdaAdj} --- \texttt{conditional nonnegative diagnostic vector}; \(\mathcal X\to\mathbb R_{+}^{K-1}\); \texttt{Observable residual energy after the cell projection and conditional Gram solve for adjacent-\allowbreak{}intensity contrasts.}
  \par\smallskip\noindent\emph{Definition.} \(\Lambda_j(X)=E_P\{[Nv_j](O)^2\mid X\}-b_{v_j}(X)^{\mathsf T}G(X)^{+}b_{v_j}(X)\), where \(N=N_B\) and \(j=1,\ldots,K-1\)
  \item \texttt{GaugAdj} --- \texttt{family of conditional augmented Gram matrices}; \(\mathcal X\to\mathbb R^{(2K^{2}+1)\times(2K^{2}+1)}\); \texttt{Observable rank-\allowbreak{}augmentation matrix testing whether an adjacent-\allowbreak{}intensity cell score adds a direction beyond the residualized measurement span.}
  \par\smallskip\noindent\emph{Definition.} For \(c_j\in\mathcal C_{\mathrm{adj}}\), \(G_j^{\mathrm{aug}}(X)=\begin{pmatrix}G(X)&b_{v_j}(X)\\b_{v_j}(X)^{\mathsf T}&E_P\{[N_Bv_j](O)^2\mid X\}\end{pmatrix}\)
  \item \texttt{Q} --- \texttt{bounded linear operator}; \(L^{2}(P)\to L^{2}(P)\); Candidate orthogonal projector onto \(\mathcal T^{\perp}\).
  \par\smallskip\noindent\emph{Definition.} \(Qf_0=N_Bf_0-h^{\mathsf T}G(X)^{+}b_{f_0}(X)\)
  \item \texttt{phieff} --- \texttt{influence-\allowbreak{}function vector}; \(L^{2}_{0}(P)^{K}\); \texttt{Candidate canonical-\allowbreak{}gradient vector.}
  \par\smallskip\noindent\emph{Definition.} \(\phi^{\mathrm{eff}}=\phi-Q\phi\)
  \item \texttt{Vraw} --- \texttt{covariance matrix}; \(\mathbb R^{K\times K}\); \texttt{Variance of the raw product-\allowbreak{}bridge gradient.}
  \par\smallskip\noindent\emph{Definition.} \(V_{\mathrm{raw}}=E_P(\phi\phi^{\mathsf T})\)
  \item \texttt{Veff} --- \texttt{covariance matrix}; \(\mathbb R^{K\times K}\); \texttt{Candidate semiparametric efficiency bound.}
  \par\smallskip\noindent\emph{Definition.} \(V_{\mathrm{eff}}=E_P(\phi^{\mathrm{eff}}\phi^{\mathrm{eff},\mathsf T})\)
  \item \texttt{eta} --- \texttt{nuisance tuple}; \texttt{Primitive and derived nuisance objects used by the score.}
  \par\smallskip\noindent\emph{Definition.} \(\eta=(m,\pi,M_A,M_Z,f,w,W,B,d,D,R,N,h,G,b_{(\cdot)},b_U)\)
  \item \texttt{tauG} --- \texttt{fixed positive real}; \((0,\varepsilon)\); \texttt{Non-\allowbreak{}oracle spectral cutoff for the fitted conditional Gram matrix.}
  \par\smallskip\noindent\emph{Definition.} \(\tau_G=\varepsilon/2\), where \(\varepsilon\) is the known lower-bound input defining the analysis model class
  \item \texttt{tauS} --- \texttt{positive training-\allowbreak{}fold-\allowbreak{}measurable sequence}; \texttt{Threshold for deleting and renormalizing low fitted measurement probabilities.}
  \par\smallskip\noindent\emph{Definition.} \(0<\tau_{S,n}<1/K\) and \(\tau_{S,n}\downarrow0\), chosen without access to the true measurement supports
  \item \texttt{Ghatplus} --- \texttt{spectral-\allowbreak{}thresholded conditional generalized inverse}; \(\mathcal X\to\mathbb R^{2K^2\times2K^2}\); \texttt{Data-\allowbreak{}driven replacement for the fitted Moore-\allowbreak{}-\allowbreak{}Penrose inverse with oracle rank matching.}
  \par\smallskip\noindent\emph{Definition.} \(\widehat G_{\tau_G}^{\dagger}\) retains and inverts exactly the eigenvalues of \(\widehat G\) that exceed \(\tau_G\)
  \item \texttt{kappaSupp} --- \texttt{nonnegative foldwise random scalar}; \texttt{Theoretical certificate that all true-\allowbreak{}zero measurement cells are deleted by support cleaning with high probability.}
  \par\smallskip\noindent\emph{Definition.} \(\kappa_{\mathrm{supp},n}=(r_{A,\infty}+r_{Z,\infty})/\tau_{S,n}\)
  \item \texttt{SigmaC} --- \texttt{contrast covariance}; \(\mathbb R^{q\times q}\); \texttt{Asymptotic covariance of the target contrasts.}
  \par\smallskip\noindent\emph{Definition.} \(\Sigma_C=CV_{\mathrm{eff}}C^{\mathsf T}\)
  \item \texttt{I0handle} --- \texttt{proof-\allowbreak{}construction handle}; \texttt{Finite-\allowbreak{}dimensional first step of the DQM lifting argument.}
  \par\smallskip\noindent\emph{Definition.} At constant \(X\), differentiate the normalized \(K\times K\times K\) feature tensor and invert the separated finite Jacobian on its observed score image to recover the centered scores of \(\pi\), \(M_A\), \(M_Z\), and \(\Gamma_Y\)
  \item \texttt{Ihandle} --- \texttt{proof-\allowbreak{}construction handle}; \texttt{Conditional extension handle for the tangent-\allowbreak{}completeness question.}
  \par\smallskip\noindent\emph{Definition.} Begin with the constant-\(X\) finite-feature inverse \(\mathfrak I_0\), choose its measurable conditional-\(X\) version, subtract the recovered \((\pi,M_A,M_Z,\Gamma_Y)\) scores, and recover the residual class-density score through the mixture equation in conditional \(L^{2}\)
  \item \texttt{dX} --- \texttt{positive integer}; \texttt{Euclidean covariate dimension in the smooth coherent-\allowbreak{}fit subclass.}
  \item \texttt{beta} --- \texttt{positive real}; \((0,\infty)\); \texttt{Common H\textbackslash{}\allowbreak{}"older smoothness index for primitive latent factors.}
  \item \texttt{rho} --- \texttt{deterministic positive sequence}; \texttt{Benchmark primitive and regular derived-\allowbreak{}operator rate on the smooth subclass;\allowbreak{} this rate alone does not control structural-\allowbreak{}zero leakage.}
  \par\smallskip\noindent\emph{Definition.} \(\rho_n=n^{-\beta/(2\beta+d_X+1)}\)
  \item \texttt{Bholder} --- \texttt{positive real}; \((0,\infty)\); \texttt{Radius of the primitive-\allowbreak{}factor H\textbackslash{}\allowbreak{}"older ball.}
  \item \texttt{MclassSmooth} --- \texttt{model subclass}; \texttt{Recognizable nonparametric subclass for the coherent sieve question.}
  \par\smallskip\noindent\emph{Definition.} \(\mathcal M^{\beta}_{K,\varepsilon}(B_{\mathrm H})\) is the compact-domain H\"older subclass of \(\mathcal M_{K,\varepsilon}\)
  \item \texttt{Sdim} --- \texttt{deterministic positive-\allowbreak{}integer sequence}; \texttt{Total tensor-\allowbreak{}product spline dimension for outcome-\allowbreak{}density estimation.}
  \par\smallskip\noindent\emph{Definition.} \(s_n\asymp n^{(d_X+1)/(2\beta+d_X+1)}\)
  \item \texttt{ell} --- \texttt{positive integer}; \(\{3,4,\ldots\}\); \texttt{Outcome-\allowbreak{}partition resolution in the Zhou-\allowbreak{}-\allowbreak{}Tchetgen Tchetgen binary approximation.}
  \item \texttt{Jbin} --- \texttt{binary latent treatment}; \(\{0,1\}\); \texttt{Binary-\allowbreak{}treatment restriction used only to state the published comparator.}
  \par\smallskip\noindent\emph{Definition.} \(J^{\mathrm{bin}}\) denotes the binary-treatment restriction of the latent state
  \item \texttt{Iztt} --- \texttt{equal-\allowbreak{}width outcome partition}; Actual discretization of \((-L,L]\) used in Supplement S8.2.
  \par\smallskip\noindent\emph{Definition.} \(\mathcal I^{\ell}=\{I_1^{\ell},\ldots,I_{\ell}^{\ell}\}\), where \(I_t^{\ell}=(-L+2(t-1)L/\ell,-L+2tL/\ell]\)
  \item \texttt{paZtt} --- \texttt{nuisance-\allowbreak{}dependent conditional cell probabilities}; \texttt{Conditional probabilities entering the Supplement S8.2 functions.}
  \par\smallskip\noindent\emph{Definition.} \(p_a^{\ell}(t,X)=P(Y\in I_t^{\ell}\mid J^{\mathrm{bin}}=a,X)\) for \(a\in\{0,1\}\)
  \item \texttt{QpartZtt} --- \texttt{partition-\allowbreak{}measurable function space}; \texttt{Piecewise-\allowbreak{}constant outcome space in Supplement S8.2.}
  \par\smallskip\noindent\emph{Definition.} \(\mathcal Q^{\ell}=\{b^{\ell}(Y,X)=\sum_{t=1}^{\ell}\mathbf 1\{Y\in I_t^{\ell}\}d_t(X):d_t\in L^2(P_X)\}\)
  \item \texttt{P1Ztt} --- \texttt{partition-\allowbreak{}dependent conditional-\allowbreak{}null space}; The \((\ell-2)\)-dimensional outcome component used by the comparator.
  \par\smallskip\noindent\emph{Definition.} \(\mathcal P_1^{\ell}=\{h\in\mathcal Q^{\ell}:E_P(h\mid J^{\mathrm{bin}},X)=0\}\)
  \item \texttt{barMA} --- \texttt{support-\allowbreak{}cleaned measurable matrix-\allowbreak{}valued function}; \(\mathcal X\to[0,1]^{K\times K}\); \texttt{Surrogate measurement matrix used by every fitted bridge and derived operator.}
  \par\smallskip\noindent\emph{Definition.} \(\bar M_A\) is obtained by thresholding and column-renormalizing the raw \(\widehat M_A\) at \(\tau_{S,n}\)
  \item \texttt{pztt} --- \texttt{family of nuisance-\allowbreak{}dependent conditional-\allowbreak{}probability functions}; The explicit Supplement S8.2 basis system for \(\mathcal P_1^{\ell}\), subject to its stated nonzero denominators.
  \par\smallskip\noindent\emph{Definition.} For \(t=2,\ldots,\ell-1\), \(p_{t-1}^{\ell}(Y,X)=\{\mathbf 1(Y\in I_1^{\ell})-p_0^{\ell}(1,X)\}/\{p_1^{\ell}(1,X)-p_0^{\ell}(1,X)\}+\{\mathbf 1(Y\in I_t^{\ell})-p_1^{\ell}(t,X)\}/\{p_0^{\ell}(t,X)-p_1^{\ell}(t,X)\}-1\)
  \item \texttt{barMZ} --- \texttt{support-\allowbreak{}cleaned measurable matrix-\allowbreak{}valued function}; \(\mathcal X\to[0,1]^{K\times K}\); \texttt{Proxy measurement matrix used by every fitted bridge and derived operator.}
  \par\smallskip\noindent\emph{Definition.} \(\bar M_Z\) is obtained by thresholding and column-renormalizing the raw \(\widehat M_Z\) at \(\tau_{S,n}\)
  \item \texttt{Bztt} --- \texttt{pair of binary measurement residual products}; \texttt{The two measurement factors used in the Supplement S8.2 orthocomplement basis.}
  \par\smallskip\noindent\emph{Definition.} \(B_0=\{A-E_P(A\mid J^{\mathrm{bin}}=0,X)\}\{Z-E_P(Z\mid J^{\mathrm{bin}}=1,X)\}\) and \(B_1=\{A-E_P(A\mid J^{\mathrm{bin}}=1,X)\}\{Z-E_P(Z\mid J^{\mathrm{bin}}=0,X)\}\)
  \item \texttt{etahat} --- \texttt{cross-\allowbreak{}fitted nuisance tuple}; \texttt{Coherent non-\allowbreak{}oracle raw-\allowbreak{}and-\allowbreak{}cleaned fitted nuisance system.}
  \par\smallskip\noindent\emph{Definition.} \(\widehat\eta=(\widehat\pi,\widehat M_A,\widehat M_Z,\widehat f,\widehat m,\bar M_A,\bar M_Z,\widehat w,\widehat W,\widehat B,\widehat d,\widehat D,\widehat R,\widehat N,\widehat h,\widehat G,\widehat G_{\tau_G}^{\dagger},\widehat b_{(\cdot)},\widehat b_U)\), where \(\widehat m_j=\int_{\mathcal Y}y\widehat f_jd\nu\), and every object after \(\bar M_A,\bar M_Z\) is derived from the same cleaned product law
  \item \texttt{Uhat} --- \texttt{estimated uncentered score vector}; \texttt{Product-\allowbreak{}bridge score evaluated at the coherent fit.}
  \par\smallskip\noindent\emph{Definition.} \(\widehat U=U(\widehat\eta)\)
  \item \texttt{Qhat} --- \texttt{estimated bounded linear operator}; \texttt{Finite projector evaluated from the support-\allowbreak{}cleaned coherent fitted factorization.}
  \par\smallskip\noindent\emph{Definition.} With \(\widehat B=\bar M_A\odot\bar M_Z\), \(\widehat N=I_{K^2}-\widehat D^{-1}\widehat B(\widehat B^{\mathsf T}\widehat D^{-1}\widehat B)^+\widehat B^{\mathsf T}\) and \(\widehat Qf_0=\widehat Nf_0-\widehat h^{\mathsf T}\widehat G_{\tau_G}^{\dagger}\widehat b_{f_0}\)
  \item \texttt{psihat} --- \texttt{estimator vector}; \(\mathbb R^{K}\); \texttt{Cross-\allowbreak{}fitted efficient estimator.}
  \par\smallskip\noindent\emph{Definition.} \(\widehat\psi=n^{-1}\sum_{i=1}^{n}\{\widehat U_{-\mathrm{fold}(i)}(O_i)-\widehat Q_{-\mathrm{fold}(i)}\widehat U_{-\mathrm{fold}(i)}(O_i)\}\)
  \item \texttt{Vhat} --- \texttt{covariance estimator}; \(\mathbb R^{K\times K}\); \texttt{Estimated efficiency bound.}
  \par\smallskip\noindent\emph{Definition.} \(\widehat V_{\mathrm{eff}}\) is the empirical covariance of the cross-fitted efficient scores centered at \(\widehat\psi\)
  \item \texttt{SigmaChat} --- \texttt{contrast covariance estimator}; \(\mathbb R^{q\times q}\); Estimated contrast covariance used for Wald and maximum-\(t\) inference.
  \par\smallskip\noindent\emph{Definition.} \(\widehat\Sigma_C=C\widehat V_{\mathrm{eff}}C^{\mathsf T}\)
  \item \texttt{VbasisZtt} --- \texttt{partition-\allowbreak{}dependent vector of functions}; \texttt{The actual finite orthocomplement basis vector in Supplement S8.2.}
  \par\smallskip\noindent\emph{Definition.} \(V^{\ell}=(B_0p_1^{\ell},B_1p_1^{\ell},\ldots,B_0p_{\ell-2}^{\ell},B_1p_{\ell-2}^{\ell})^{\mathsf T}\)
  \item \texttt{Qztt} --- \texttt{family of bounded linear operators}; \(L^{2}(P)\to L^{2}(P)\); \texttt{Supplement S8.2 finite-\allowbreak{}partition projection in the binary restriction.}
  \par\smallskip\noindent\emph{Definition.} \(Q_{\ell}^{\mathrm{ZTT}}f=E_P(fV^{\ell,\mathsf T})\{E_P(V^{\ell}V^{\ell,\mathsf T})\}^{-1}V^{\ell}\)
  \item \texttt{Fithandle} --- \texttt{estimator-\allowbreak{}construction handle}; \texttt{Non-\allowbreak{}oracle support-\allowbreak{}cleaned starting move for coherent nuisance construction.}
  \par\smallskip\noindent\emph{Definition.} On \(\mathcal M^{\beta}_{K,\varepsilon}(B_{\mathrm H})\), fit one raw label-aligned tensor-product B-spline product likelihood, derive \(\widehat m\) from \(\widehat f\), threshold and renormalize \(\widehat M_A,\widehat M_Z\) into \(\bar M_A,\bar M_Z\), derive every bridge and operator from the cleaned law, and use \(\tau_G=\varepsilon/2\); the likelihood-curvature, localization, boundary, and uniform-rate proof remains the explicit open part
  \item \texttt{Vztt} --- \texttt{family of covariance matrices}; \(\mathbb R^{2\times2}\); \texttt{Approximate variance delivered by the published outcome-\allowbreak{}partition projection.}
  \par\smallskip\noindent\emph{Definition.} \(V_{\ell}^{\mathrm{ZTT}}=E_P[(\phi-Q_{\ell}^{\mathrm{ZTT}}\phi)(\phi-Q_{\ell}^{\mathrm{ZTT}}\phi)^{\mathsf T}]\) in the binary restriction
  \item \texttt{rAc} --- \texttt{nonnegative foldwise random scalar}; \texttt{Cleaned surrogate measurement-\allowbreak{}rate envelope.}
  \par\smallskip\noindent\emph{Definition.} \(r_A^c=r_A+\tau_{S,n}\)
  \item \texttt{rZc} --- \texttt{nonnegative foldwise random scalar}; \texttt{Cleaned proxy measurement-\allowbreak{}rate envelope.}
  \par\smallskip\noindent\emph{Definition.} \(r_Z^c=r_Z+\tau_{S,n}\)
  \item \texttt{reta} --- \texttt{nonnegative rate vector}; \texttt{Foldwise raw,\allowbreak{} cleaned,\allowbreak{} derived-\allowbreak{}operator,\allowbreak{} spectral,\allowbreak{} and support-\allowbreak{}screening errors.}
  \par\smallskip\noindent\emph{Definition.} \(r_\eta=(r_m,r_\pi,r_A,r_Z,r_A^c,r_Z^c,r_f,r_w,r_W,r_B,r_d,r_D,r_R,r_N,r_h,r_G,r_b,r_{A,\infty},r_{Z,\infty},r_{G,\infty},\kappa_{\mathrm{supp},n})\); \(r_A,r_Z\) are raw conditional \(L^2(P)\) errors, \(r_A^c,r_Z^c\) are cleaned envelopes, \(r_B,\ldots,r_b\) are cleaned-law plug-in errors, and the essential-supremum measurement errors are screening quantities only
  \item \texttt{Rsum} --- \texttt{nonnegative random scalar}; \texttt{Aggregate cleaned-\allowbreak{}law score-\allowbreak{}stability and Gram-\allowbreak{}selection error.}
  \par\smallskip\noindent\emph{Definition.} \(R_{\Sigma}=r_m+r_\pi+r_A^c+r_Z^c+r_f+r_w+r_W+r_B+r_d+r_D+r_R+r_N+r_h+r_G+r_b+r_{G,\infty}\)
  \item \texttt{Delta} --- \texttt{nonnegative random scalar}; \texttt{Coherent cleaned-\allowbreak{}fit quadratic mean-\allowbreak{}bias bound.}
  \par\smallskip\noindent\emph{Definition.} \(\Delta_n=r_m(r_\pi+r_A^c+r_Z^c)+r_A^cr_Z^c+(r_A^c+r_Z^c)R_{\Sigma}\)
\end{itemize}

\section{Assumptions}

\begin{assumption}[ass:iid]\label{ass:iid}
\(O_1,\ldots,O_n\stackrel{\mathrm{iid}}{\sim}P\).
\par\smallskip\noindent\emph{Standard} (Independent sampling, \cite{ZhouTchetgenTchetgen2026}).
\end{assumption}

\begin{assumption}[ass:fixed-k]\label{ass:fixed-k}
\(K\geq3\) is fixed as \(n\to\infty\).
\par\smallskip\noindent\emph{Standard} (Fixed finite latent cardinality, \cite{AllmanMatiasRhodes2009}).
\end{assumption}

\begin{assumption}[ass:bounded-outcome]\label{ass:bounded-outcome}
\(|Y|\leq L\) almost surely.
\par\smallskip\noindent\emph{Standard} (Bounded outcome, \cite{ZhouTchetgenTchetgen2026}).
\end{assumption}

\begin{assumption}[ass:dominated-outcome]\label{ass:dominated-outcome}
For every \(j\), the conditional law of \(Y\) given \((J=j,X=x)\) has density \(f_j(\cdot\mid x)\) with respect to \(\nu\) on \(\mathcal Y\).
\par\smallskip\noindent\emph{Standard} (Dominated conditional outcome model, \cite{GuoOgburnShpitser2026}).
\end{assumption}

\begin{assumption}[ass:common-support]\label{ass:common-support}
\(f_j(y\mid x)>0\) for every \(j\), for \(\nu\)-almost every \(y\in\mathcal Y\), and for \(P_X\)-almost every \(x\).
\par\smallskip\noindent\emph{Standard} (Common-support variant, \cite{GuoOgburnShpitser2026}).
\end{assumption}

\begin{assumption}[ass:bounded-density]\label{ass:bounded-density}
\(\sup_{j,y,x}f_j(y\mid x)\leq\varepsilon^{-1}\).
\par\smallskip\noindent\emph{Novel.} This proposal-specific envelope is plausible after restricting birth weight to a prespecified clinically meaningful compact range and excluding arbitrarily concentrated within-bin latent-class laws; smooth location-scale families with scales bounded away from zero and small common-support perturbations satisfy it.
\end{assumption}

\begin{assumption}[ass:three-view-independence]\label{ass:three-view-independence}
\(Y\), \(A\), and \(Z\) are mutually conditionally independent given \((J,X)\).
\par\smallskip\noindent\emph{Standard} (Three-view conditional independence, \cite{AllmanMatiasRhodes2009}).
\end{assumption}

\begin{assumption}[ass:consistency]\label{ass:consistency}
\(Y=Y(J)\) almost surely.
\par\smallskip\noindent\emph{Standard} (Consistency, \cite{ZhouTchetgenTchetgen2026}).
\end{assumption}

\begin{assumption}[ass:latent-ignorability]\label{ass:latent-ignorability}
\(Y(j)\perp J\mid X\) for every \(j\in\mathcal J\).
\par\smallskip\noindent\emph{Standard} (Latent treatment ignorability, \cite{ZhouTchetgenTchetgen2026}).
\end{assumption}

\begin{assumption}[ass:latent-positivity]\label{ass:latent-positivity}
\(\min_{j}\pi_j(X)\geq\varepsilon\) almost surely.
\par\smallskip\noindent\emph{Standard} (Uniform latent-treatment positivity, \cite{ZhouTchetgenTchetgen2026}).
\end{assumption}

\begin{assumption}[ass:cell-positivity]\label{ass:cell-positivity}
\(\min_{a,z}d_{az}(Y,X)\geq\varepsilon\) almost surely.
\par\smallskip\noindent\emph{Standard} (Uniform observed-cell positivity, \cite{ZhouTchetgenTchetgen2026}).
\end{assumption}

\begin{assumption}[ass:surrogate-rank]\label{ass:surrogate-rank}
\(\sigma_{\min}\{M_A(X)\}\geq\varepsilon\) almost surely.
\par\smallskip\noindent\emph{Novel.} This proposal-specific quantitative separation strengthens the qualitative nonsingularity used by Bian, Zhou, and Cui. In a six-bin maternal-smoking design it is plausible when the validation confusion matrix for questionnaire responses has six distinguishable columns, so no two true intensity bins induce nearly identical self-report distributions; the smallest singular value is directly diagnosable in the fitted measurement design.
\end{assumption}

\begin{assumption}[ass:proxy-rank]\label{ass:proxy-rank}
\(\sigma_{\min}\{M_Z(X)\}\geq\varepsilon\) almost surely.
\par\smallskip\noindent\emph{Novel.} This proposal-specific quantitative separation strengthens qualitative proxy nonsingularity. It is plausible when the six cotinine categories are cut at assay-calibrated thresholds whose class-conditional response profiles remain distinct across the six latent smoking intensities; a paired validation sample can screen the minimum singular value before analysis.
\end{assumption}

\begin{assumption}[ass:outcome-feature-rank]\label{ass:outcome-feature-rank}
\(\sigma_{\min}\{\Gamma_Y(X)\}\geq\varepsilon\) almost surely.
\par\smallskip\noindent\emph{Novel.} This proposal-specific margin requires the six latent exposure levels to induce quantitatively distinct birth-weight feature profiles, beyond qualitative factor rank. Clinically prespecified birth-weight intervals or low-order spline features can satisfy it when increasing smoking intensity shifts more than a single common mean and the six feature vectors are not nearly collinear.
\end{assumption}

\begin{assumption}[ass:modal-label-margin]\label{ass:modal-label-margin}
\(M_A(X)_{jj}-\max_{a\neq j}M_A(X)_{aj}\geq\varepsilon\) for every \(j\) almost surely.
\par\smallskip\noindent\emph{Novel.} This proposal-specific uniform margin strengthens the strict modal alignment used by Bian, Zhou, and Cui. For six maternal-smoking bins it is plausible when respondents are more likely to report their own intensity bin than any particular alternative by a stable amount after covariate adjustment; repeated questionnaire or validation data can estimate the modal gaps.
\end{assumption}

\begin{assumption}[ass:posterior-gram-gap]\label{ass:posterior-gram-gap}
The smallest eigenvalue of \(R(Y,X)^{\mathsf T}D(Y,X)R(Y,X)\) is at least \(\varepsilon\) almost surely.
\par\smallskip\noindent\emph{Novel.} This proposal-specific margin requires the six posterior class profiles across the \(36\) self-report--cotinine cells to retain independent weighted variation at each birth-weight and covariate stratum. It is plausible when the two instruments make complementary errors and the augmented study deliberately populates every paired category through assay calibration and oversampling of discordant reports; the recovered \(R^{\mathsf T}DR\) eigenvalues provide a design diagnostic.
\end{assumption}

\begin{assumption}[ass:projection-gram-rank]\label{ass:projection-gram-rank}
\(\operatorname{rank}\{G(X)\}\) is constant almost surely.
\par\smallskip\noindent\emph{Novel.} This proposal-specific regular-stratum condition excludes covariate strata at which residualized questionnaire or cotinine score directions appear or disappear exactly. In the six-bin augmented study it is plausible under a fixed measurement protocol and common category support, with sufficiently populated \(36\)-cell tables so the same estimable linear combinations of the \(72\) raw measurement directions persist across covariates.
\end{assumption}

\begin{assumption}[ass:projection-gram-gap]\label{ass:projection-gram-gap}
Every nonzero eigenvalue of \(G(X)\) is at least \(\varepsilon\) almost surely.
\par\smallskip\noindent\emph{Novel.} This proposal-specific gap prevents the stable residualized measurement directions from becoming nearly collinear. In the six-bin maternal-smoking design it is plausible when self-report and cotinine retain complementary conditional variation after posterior outcome adjustment and the study enriches sparse discordant cells; the nonzero eigenvalues of the recovered conditional Gram matrix are estimable design-strength diagnostics.
\end{assumption}

\begin{assumption}[ass:smooth-domain]\label{ass:smooth-domain}
\(X\in[0,1]^{d_X}\) almost surely.
\par\smallskip\noindent\emph{Standard} (Compact Euclidean covariate support, \cite{Tsybakov2009}).
\end{assumption}

\begin{assumption}[ass:holder-factors]\label{ass:holder-factors}
\((\pi,M_A,M_Z,(f_j)_{j=1}^{K})\in[C^{\beta}_{B_{\mathrm H}}([0,1]^{d_X})]^{K+2K^{2}}\times[C^{\beta}_{B_{\mathrm H}}(\mathcal Y\times[0,1]^{d_X})]^{K}\).
\par\smallskip\noindent\emph{Standard} (H\"older-smooth nonparametric factors, \cite{Tsybakov2009}).
\end{assumption}

\begin{assumption}[ass:contrast-rank]\label{ass:contrast-rank}
\(\operatorname{rank}(C)=q\).
\par\smallskip\noindent\emph{Standard} (Full-row-rank contrast, \cite{ZhouTchetgenTchetgen2026}).
\end{assumption}

\begin{assumption}[ass:contrast-centering]\label{ass:contrast-centering}
\(C\mathbf 1_K=0\).
\par\smallskip\noindent\emph{Standard} (Treatment-contrast normalization, \cite{BodoryBusshoffLechner2022}).
\end{assumption}

\begin{assumption}[ass:contrast-information]\label{ass:contrast-information}
\(\Sigma_C\) is positive definite.
\par\smallskip\noindent\emph{Standard} (Nonsingular target information, \cite{TsiatisMa2004}).
\end{assumption}

\section{Definitions}

\begin{definition}[def:model-class]\label{def:model-class}
\par\smallskip\noindent\emph{Defined object.} \texttt{Regular hidden ordered-\allowbreak{}treatment model}
\par\smallskip\noindent\emph{Construction.} \(\mathcal M_{K,\varepsilon}=\{P:\ P\text{ obeys the listed member properties for fixed }K\}\)
\par\smallskip\noindent\emph{Member properties.} \texttt{ass:\allowbreak{}fixed-\allowbreak{}k}; \texttt{ass:\allowbreak{}bounded-\allowbreak{}outcome}; \texttt{ass:\allowbreak{}dominated-\allowbreak{}outcome}; \texttt{ass:\allowbreak{}common-\allowbreak{}support}; \texttt{ass:\allowbreak{}bounded-\allowbreak{}density}; \texttt{ass:\allowbreak{}three-\allowbreak{}view-\allowbreak{}independence}; \texttt{ass:\allowbreak{}consistency}; \texttt{ass:\allowbreak{}latent-\allowbreak{}ignorability}; \texttt{ass:\allowbreak{}latent-\allowbreak{}positivity}; \texttt{ass:\allowbreak{}cell-\allowbreak{}positivity}; \texttt{ass:\allowbreak{}surrogate-\allowbreak{}rank}; \texttt{ass:\allowbreak{}proxy-\allowbreak{}rank}; \texttt{ass:\allowbreak{}outcome-\allowbreak{}feature-\allowbreak{}rank}; \texttt{ass:\allowbreak{}modal-\allowbreak{}label-\allowbreak{}margin}; \texttt{ass:\allowbreak{}posterior-\allowbreak{}gram-\allowbreak{}gap}; \texttt{ass:\allowbreak{}projection-\allowbreak{}gram-\allowbreak{}rank}; \texttt{ass:\allowbreak{}projection-\allowbreak{}gram-\allowbreak{}gap}.
\end{definition}

\begin{definition}[def:continuous-witness]\label{def:continuous-witness}
\par\smallskip\noindent\emph{Defined object.} \texttt{Positive three-\allowbreak{}class continuous witness}
\par\smallskip\noindent\emph{Construction.} \(P^{\star}\) is induced by \(X\equiv0\), \(K=L=3\), \(\mathcal Y=[0,3]\), \(\pi=(0.3,0.4,0.3)^{\mathsf T}\), \(M_A=0.7I_3+0.1\mathbf 1_3\mathbf 1_3^{\mathsf T}\), \(M_Z=0.4I_3+0.2\mathbf 1_3\mathbf 1_3^{\mathsf T}\), and \(f_j(y)=\sum_{b=0}^{2}H_{bj}\mathbf 1\{b\leq y<b+1\}\), where \(H=\begin{pmatrix}0.7&0.2&0.1\\0.2&0.6&0.2\\0.1&0.2&0.7\end{pmatrix}\), \(\nu\) is Lebesgue measure, \(h_0(y)=(\mathbf 1\{0\leq y<1\},\mathbf 1\{1\leq y<2\},\mathbf 1\{2\leq y\leq3\})^{\mathsf T}\), \(Y^{\mathrm{pot}}\perp J\) with marginal density \(f_j\) for \(Y(j)\), \(Y=Y(J)\), and \(A\), \(Z\), and \(Y^{\mathrm{pot}}\) are mutually conditionally independent given \(J\)
\par\smallskip\noindent\emph{Inputs.} \texttt{pi}; \texttt{MA}; \texttt{MZ}; \texttt{f}; \texttt{nu}; \texttt{h0}; \texttt{Ypot}.
\end{definition}

\begin{definition}[def:raw-product-bridge]\label{def:raw-product-bridge}
\par\smallskip\noindent\emph{Defined object.} \texttt{Raw product-\allowbreak{}bridge gradient}
\par\smallskip\noindent\emph{Construction.} \(U_j=m_j(X)+a_j(A,X)z_j(Z,X)\{Y-m_j(X)\}/\pi_j(X)\), with \(\phi_j=U_j-\psi_j\)
\par\smallskip\noindent\emph{Inputs.} \texttt{m}; \texttt{pi}; \texttt{bridgeA}; \texttt{bridgeZ}.
\end{definition}

\begin{definition}[def:latent-score-image]\label{def:latent-score-image}
\par\smallskip\noindent\emph{Defined object.} \texttt{Latent additive score image}
\par\smallskip\noindent\emph{Construction.} \(\overline{\mathcal S_{\mathrm{lat}}}^{\,L^{2}(P)}\), where \(\mathcal S_{\mathrm{lat}}\) is the observed conditional-expectation image of centered factor scores
\par\smallskip\noindent\emph{Inputs.} \texttt{Slat}.
\end{definition}

\begin{definition}[def:finite-projector]\label{def:finite-projector}
\par\smallskip\noindent\emph{Defined object.} \texttt{Pointwise-\allowbreak{}cell and conditional-\allowbreak{}Gram projector}
\par\smallskip\noindent\emph{Construction.} Set \(W=\operatorname{diag}(w)\), \(D=\operatorname{diag}(Bw)\), and \(R=D^{-1}BW\). Define the first elimination by \(N=N_B=I_{K^{2}}-D^{-1}B(B^{\mathsf T}D^{-1}B)^{+}B^{\mathsf T}=I_{K^{2}}-R(R^{\mathsf T}DR)^{+}R^{\mathsf T}D\), residualize the \(2K^{2}\) measurement vectors as \(h=N_Bc\), set \(G(X)=E_P(hh^{\mathsf T}\mid X)\) and \(b_{f_0}(X)=E_P\{h(N_Bf_0)\mid X\}\), and define \(Qf_0=N_Bf_0-h^{\mathsf T}G(X)^{+}b_{f_0}(X)\)
\par\smallskip\noindent\emph{Inputs.} \texttt{w}; \texttt{W}; \texttt{B}; \texttt{R}; \texttt{D}; \texttt{cellc}.
\end{definition}

\begin{definition}[def:k3-projector-audit]\label{def:k3-projector-audit}
\par\smallskip\noindent\emph{Defined object.} Exact symbolic \(K=3\) projector audit
\par\smallskip\noindent\emph{Construction.} In \(\mathfrak A_3\), use \(W=\operatorname{diag}(w_1,w_2,w_3)\) with positive nonconstant weights, \(D=\operatorname{diag}(Bw)\), \(R=D^{-1}BW\), and \(N_B=I_9-D^{-1}B(B^{\mathsf T}D^{-1}B)^{-1}B^{\mathsf T}\); verify the six matrix identities in Lemma 3A by exact symbolic multiplication and range--kernel algebra.
\par\smallskip\noindent\emph{Inputs.} \texttt{K3audit}.
\end{definition}

\begin{definition}[def:canonical-gradient]\label{def:canonical-gradient}
\par\smallskip\noindent\emph{Defined object.} \texttt{Projected gradient}
\par\smallskip\noindent\emph{Construction.} \(\phi^{\mathrm{eff}}=\phi-Q\phi\) and \(V_{\mathrm{eff}}=E_P(\phi^{\mathrm{eff}}\phi^{\mathrm{eff},\mathsf T})\)
\par\smallskip\noindent\emph{Inputs.} \texttt{phi}; \texttt{Q}.
\end{definition}

\begin{definition}[def:observable-gain-diagnostic]\label{def:observable-gain-diagnostic}
\par\smallskip\noindent\emph{Defined object.} \texttt{Adjacent-\allowbreak{}intensity efficiency-\allowbreak{}gain diagnostic}
\par\smallskip\noindent\emph{Construction.} For each \(c_j\in\mathcal C_{\mathrm{adj}}\), form the recovered cell vector \(v_j\), apply the measurement-form projector \(N_B=I_{K^{2}}-D^{-1}B(B^{\mathsf T}D^{-1}B)^{+}B^{\mathsf T}\), augment \(G(X)\) to \(G_j^{\mathrm{aug}}(X)=\begin{pmatrix}G(X)&b_{v_j}(X)\\b_{v_j}(X)^{\mathsf T}&E_P\{[N_Bv_j](O)^2\mid X\}\end{pmatrix}\), and take its Schur residual \(\Lambda_j(X)=E_P\{[N_Bv_j](O)^2\mid X\}-b_{v_j}(X)^{\mathsf T}G(X)^{+}b_{v_j}(X)\)
\par\smallskip\noindent\emph{Inputs.} \texttt{vAdj}; \texttt{N}; \texttt{G}; \texttt{bf}.
\end{definition}

\begin{definition}[def:coherent-fit]\label{def:coherent-fit}
\par\smallskip\noindent\emph{Defined object.} \texttt{Support-\allowbreak{}cleaned coherent cross-\allowbreak{}fitted nuisance system}
\par\smallskip\noindent\emph{Construction.} For each training fold, first fit one label-aligned raw product factorization \((\widehat\pi,\widehat M_A,\widehat M_Z,\widehat f)\) on the declared support \(\mathcal Y\), without using an unknown measurement support, the true conditional Gram matrix, or its rank.  Define
\[
 \widehat m_j(x)=\int_{\mathcal Y}y\widehat f_j(y\mid x)d\nu(y).
\]
Choose a positive training-fold-measurable \(\tau_{S,n}<1/K\), and apply the support-cleaning definition columnwise to obtain \(\bar M_A\) and \(\bar M_Z\).  Every bridge and every subsequent derived object uses the cleaned matrices, never the raw matrices.  In particular, set
\[
 \widehat B=\bar M_A\odot\bar M_Z,
 \quad \widehat W=\operatorname{diag}(\widehat w),
 \quad \widehat d=\widehat B\widehat w,
 \quad \widehat D=\operatorname{diag}(\widehat d),
 \quad \widehat R=\widehat D^{-1}\widehat B\widehat W,
\]
and
\[
 \widehat N=I_{K^2}-\widehat D^{-1}\widehat B
 (\widehat B^{\mathsf T}\widehat D^{-1}\widehat B)^+
 \widehat B^{\mathsf T},
 \qquad \widehat h=\widehat N\widehat c.
\]
Form \(\widehat G,\widehat b_{(\cdot)},\widehat b_U\) from this same cleaned product law.  Retain only fits satisfying the computable half-margin restrictions
\[
 \min_j\widehat\pi_j\geq\varepsilon/2,
 \quad \min_{a,z}\widehat d_{az}\geq\varepsilon/2,
 \quad \sigma_{\min}(\bar M_A)\wedge\sigma_{\min}(\bar M_Z)\geq\varepsilon/2,
 \quad \lambda_{\min}(\widehat R^{\mathsf T}\widehat D\widehat R)\geq\varepsilon/2.
\]
With the known cutoff \(\tau_G=\varepsilon/2\), form \(\widehat G_{\tau_G}^{\dagger}\) from the fitted spectrum and define
\[
 \widehat Qf_0=\widehat Nf_0-
 \widehat h^{\mathsf T}\widehat G_{\tau_G}^{\dagger}\widehat b_{f_0}.
\]
The raw essential-supremum measurement errors are used only to certify the high-probability cleaning event; no true zero pattern is queried, and no beta-min condition is imposed on positive measurement probabilities.
\par\smallskip\noindent\emph{Inputs.} \texttt{Fithandle}; \texttt{tauG}; \texttt{tauS}; \texttt{barMA}; \texttt{barMZ}.
\end{definition}

\begin{definition}[def:quadratic-remainder]\label{def:quadratic-remainder}
\par\smallskip\noindent\emph{Defined object.} \texttt{Support-\allowbreak{}cleaned coherent nuisance remainder}
\par\smallskip\noindent\emph{Construction.} Let
\[
 \kappa_{\mathrm{supp},n}
 =\frac{r_{A,\infty}+r_{Z,\infty}}{\tau_{S,n}},
 \qquad
 r_A^c=r_A+\tau_{S,n},
 \qquad
 r_Z^c=r_Z+\tau_{S,n}.
\]
Here \(r_A,r_Z\) are the conditional \(L^2(P)\) errors of the raw measurement fits, while \(r_A^c,r_Z^c\) are cleaned-rate envelopes.  Every coordinate \(r_B,r_d,r_D,r_R,r_N,r_h,r_G,r_b\) is the error of the plug-in computed from \((\bar M_A,\bar M_Z)\).  Set
\[
\begin{aligned}
 R_\Sigma={}&r_m+r_\pi+r_A^c+r_Z^c+r_f+r_w+r_W+r_B+r_d+r_D
 +r_R+r_N+r_h+r_G+r_b+r_{G,\infty},\\
 \Delta_n={}&r_m(r_\pi+r_A^c+r_Z^c)+r_A^cr_Z^c
 +(r_A^c+r_Z^c)R_\Sigma.
\end{aligned}
\]
The raw errors \(r_{A,\infty},r_{Z,\infty}\) occur only in \(\kappa_{\mathrm{supp},n}\), which certifies exact removal of true-zero cells.  There is no additive first-order leakage term after cleaning.
\par\smallskip\noindent\emph{Inputs.} \texttt{reta}; \texttt{tauS}; \texttt{rAc}; \texttt{rZc}; \texttt{kappaSupp}.
\end{definition}

\begin{definition}[def:cross-fitted-estimator]\label{def:cross-fitted-estimator}
\par\smallskip\noindent\emph{Defined object.} \texttt{Efficient cross-\allowbreak{}fitted estimator}
\par\smallskip\noindent\emph{Construction.} \(\widehat\psi=n^{-1}\sum_{i=1}^{n}\{\widehat U_{-\mathrm{fold}(i)}(O_i)-\widehat Q_{-\mathrm{fold}(i)}\widehat U_{-\mathrm{fold}(i)}(O_i)\}\)
\par\smallskip\noindent\emph{Inputs.} \texttt{etahat}; \texttt{Uhat}; \texttt{Qhat}.
\end{definition}

\begin{definition}[def:constant-x-inverse-handle]\label{def:constant-x-inverse-handle}
\par\smallskip\noindent\emph{Defined object.} Constant-\(X\) finite-feature inverse handle
\par\smallskip\noindent\emph{Construction.} Use \(\mathfrak I_0\): when \(X\) is degenerate, augment the \(K\times K\times K\) feature tensor \(T_u=E\{h_{0,u}(Y)e_Ae_Z^{\mathsf T}\}\) by the observed normalization slice \(S=P(A,Z)\), differentiate \((S,T_1,\ldots,T_K)\), and invert its separated finite Jacobian on the observed score image to recover centered scores for \(\pi\), \(M_A\), \(M_Z\), and \(\Gamma_Y\).
\par\smallskip\noindent\emph{Inputs.} \texttt{pi}; \texttt{MA}; \texttt{MZ}; \texttt{GammaY}.
\end{definition}

\begin{definition}[def:differential-inverse-handle]\label{def:differential-inverse-handle}
\par\smallskip\noindent\emph{Defined object.} \texttt{Differential factor-\allowbreak{}lifting handle}
\par\smallskip\noindent\emph{Construction.} Use \(\mathfrak I\): start from \(\mathfrak I_0\), select a measurable conditional-\(X\) finite-factor inverse with a uniform operator bound, subtract the recovered finite factor scores, and invert the remaining mixture equation for each support-preserving outcome-density score in conditional \(L^{2}\)
\par\smallskip\noindent\emph{Inputs.} \texttt{GammaY}; \texttt{pi}; \texttt{MA}; \texttt{MZ}; \texttt{f}.
\end{definition}

\begin{definition}[def:smooth-subclass]\label{def:smooth-subclass}
\par\smallskip\noindent\emph{Defined object.} \texttt{H\textbackslash{}\allowbreak{}"older coherent-\allowbreak{}fit subclass}
\par\smallskip\noindent\emph{Construction.} \(\mathcal M^{\beta}_{K,\varepsilon}(B_{\mathrm H})=\{P\in\mathcal M_{K,\varepsilon}:P\text{ obeys the compact-domain and H\"older-factor member properties for }d_X,\beta,B_{\mathrm H}\}\)
\par\smallskip\noindent\emph{Member properties.} \texttt{ass:\allowbreak{}smooth-\allowbreak{}domain}; \texttt{ass:\allowbreak{}holder-\allowbreak{}factors}.
\end{definition}

\begin{definition}[def:ztt-partition-approximation]\label{def:ztt-partition-approximation}
\par\smallskip\noindent\emph{Defined object.} \texttt{Zhou-\allowbreak{}-\allowbreak{}Tchetgen Tchetgen finite-\allowbreak{}partition variance sequence}
\par\smallskip\noindent\emph{Construction.} In the binary restriction, partition \((-L,L]\) into \(\mathcal I^{\ell}\), form \(\mathcal P_1^{\ell}=\{h\in\mathcal Q^{\ell}:E_P(h\mid J^{\mathrm{bin}},X)=0\}\), use the Supplement S8.2 conditional-probability basis \((p_1^{\ell},\ldots,p_{\ell-2}^{\ell})\), set \(V^{\ell}=(B_0p_1^{\ell},B_1p_1^{\ell},\ldots,B_0p_{\ell-2}^{\ell},B_1p_{\ell-2}^{\ell})^{\mathsf T}\), and project by \(Q_{\ell}^{\mathrm{ZTT}}f=E_P(fV^{\ell,\mathsf T})\{E_P(V^{\ell}V^{\ell,\mathsf T})\}^{-1}V^{\ell}\)
\par\smallskip\noindent\emph{Inputs.} \texttt{Iztt}; \texttt{paZtt}; \texttt{pztt}; \texttt{Bztt}; \texttt{VbasisZtt}; \texttt{phi}.
\end{definition}

\begin{definition}[def:coherent-fit-handle]\label{def:coherent-fit-handle}
\par\smallskip\noindent\emph{Defined object.} \texttt{Non-\allowbreak{}oracle support-\allowbreak{}cleaned nuisance-\allowbreak{}rate handle}
\par\smallskip\noindent\emph{Construction.} Use \(\mathfrak F\) as follows.  On \(\mathcal M^{\beta}_{K,\varepsilon}(B_{\mathrm H})\), fit one raw tensor-product B-spline product likelihood for \((\pi,M_A,M_Z,f)\), with \(s_n\asymp n^{(d_X+1)/(2\beta+d_X+1)}\), and define \(\widehat m_j=\int_{\mathcal Y}y\widehat f_jd\nu\).  Among the \(K!\) common column permutations, choose the maximizer of the empirical training-sample minimum diagonal-minus-largest-off-diagonal self-report probability.  Choose \(\tau_{S,n}\) from the fold size and training data only, threshold and renormalize the columns of the raw \(\widehat M_A,\widehat M_Z\) to obtain \(\bar M_A,\bar M_Z\), and retain only the computable cleaned-law half-margin restrictions in the coherent-fit definition.  Derive \(\widehat w,\widehat W,\widehat B,\widehat d,\widehat D,\widehat R,\widehat N,\widehat h,\widehat G,\widehat b_{(\cdot)},\widehat b_U\), and \(\widehat Q\) from \((\widehat\pi,\bar M_A,\bar M_Z,\widehat f)\).  Use \(\tau_G=\varepsilon/2\) for the fitted spectral inverse.  Neither a true support, the true \(G\), nor a true rank is an input.  The unresolved part of this handle is to supply the spline boundary convention, constrained likelihood curvature, localization, and uniform-rate facts requested in the coherent-fit-rate question.
\par\smallskip\noindent\emph{Inputs.} \texttt{MclassSmooth}; \texttt{Sdim}; \texttt{pi}; \texttt{MA}; \texttt{MZ}; \texttt{barMA}; \texttt{barMZ}; \texttt{f}; \texttt{tauG}; \texttt{tauS}.
\end{definition}

\begin{definition}[def:spectral-thresholded-gram-inverse]\label{def:spectral-thresholded-gram-inverse}
\par\smallskip\noindent\emph{Defined object.} \texttt{Data-\allowbreak{}driven conditional Gram inverse}
\par\smallskip\noindent\emph{Construction.} For a fitted conditional Gram matrix with spectral decomposition
\[
 \widehat G(x)=\sum_{s=1}^{2K^2}\widehat\lambda_s(x)
 \widehat v_s(x)\widehat v_s(x)^{\mathsf T},
 \qquad \widehat\lambda_1(x)\geq\cdots\geq\widehat\lambda_{2K^2}(x)\geq0,
\]
set \(\tau_G=\varepsilon/2\) and define the Borel spectral functional
\[
 \widehat G_{\tau_G}^{\dagger}(x)
 =\sum_{s:\,\widehat\lambda_s(x)>\tau_G}
 \widehat\lambda_s(x)^{-1}
 \widehat v_s(x)\widehat v_s(x)^{\mathsf T}.
\]
Thus the fitted rank is the observable count \(\widehat r_G(x)=\sum_s\mathbf 1\{\widehat\lambda_s(x)>\tau_G\}\); neither \(G(x)\) nor its rank is an input.
\par\smallskip\noindent\emph{Inputs.} \texttt{etahat}; \texttt{tauG}.
\end{definition}

\begin{definition}[def:support-leakage-diagnostic]\label{def:support-leakage-diagnostic}
\par\smallskip\noindent\emph{Defined object.} \texttt{Threshold-\allowbreak{}and-\allowbreak{}renormalize measurement-\allowbreak{}support cleaning}
\par\smallskip\noindent\emph{Construction.} Fix a positive training-fold-measurable threshold \(\tau_{S,n}\downarrow0\) with \(\tau_{S,n}<1/K\), chosen without the true measurement supports.  From the raw fitted measurement matrices define, columnwise,
\[
 \bar M_{A,aj}(x)=
 \frac{\widehat M_A(x)_{aj}\mathbf 1\{\widehat M_A(x)_{aj}>\tau_{S,n}\}}
 {\sum_{b=1}^K\widehat M_A(x)_{bj}\mathbf 1\{\widehat M_A(x)_{bj}>\tau_{S,n}\}},
\]
and define \(\bar M_Z\) analogously.  The denominator is positive because a probability column has at least one entry not smaller than \(1/K>\tau_{S,n}\).  These cleaned matrices, rather than \(\widehat M_A,\widehat M_Z\), generate every fitted bridge, \(\widehat B,\widehat d,\widehat D,\widehat R,\widehat N,\widehat h,\widehat G,\widehat b_{(\cdot)}\), and \(\widehat Q\).  The algorithm does not query a true support.  If support sets are supplied as observed design information, their zeros may additionally be imposed, but this is not required.
\par\smallskip\noindent\emph{Inputs.} \texttt{etahat}; \texttt{tauS}; \texttt{MA}; \texttt{MZ}.
\end{definition}

\section{Main results}

\begin{theorem}[thm:1-identification]\label{thm:1-identification}
For every \(P\in\mathcal M_{K,\varepsilon}\), the observable augmented moment system \(S(x)=\{P(A=a,Z=z\mid X=x)\}_{a,z}\) and \(T_u(x)=E_P\{h_{0,u}(Y)e_Ae_Z^{\mathsf T}\mid X=x\}\), \(u=1,\ldots,K\), identifies \((\pi,M_A,M_Z,\Gamma_Y)\) up to a common permutation; together with the full observed conditional density \(p_A(y\mid x)\), it identifies \(f(y\mid x)=\Pi^{-1}M_A^{-1}p_A(y\mid x)\). The modal margin selects the unique ordered labeling, and \(\psi_j=E_P\{m_j(X)\}=E\{Y(j)\}\) for every \(j\in\mathcal J\).
\end{theorem}
\begin{proof}
Fix measurable versions of all conditional laws and fix \(x\) outside the common \(P_X\)-null set on which the assumptions may fail.  Abbreviate
\[
A_0=M_A(x),\qquad Z_0=M_Z(x),\qquad \Pi=\operatorname{diag}\{\pi(x)\},
\qquad \gamma_j=\Gamma_Y(x)_{\cdot j}.
\]
Conditional three-view independence in \(\mathcal M_{K,\varepsilon}\) gives, for the observable matrices in the statement,
\[
S=A_0\Pi Z_0^{\mathsf T},\qquad
T_u=A_0\Pi\operatorname{diag}(\gamma_{u1},\ldots,\gamma_{uK})Z_0^{\mathsf T},
\quad u=1,\ldots,K. \tag{1}
\]
The latent-positivity, surrogate-rank, proxy-rank, and outcome-feature-rank assumptions make \(\Pi,A_0,Z_0,\Gamma_Y(x)\) nonsingular.  Thus \(S\) is nonsingular, and (1) implies
\[
T_uS^{-1}=A_0\operatorname{diag}(\gamma_{u1},\ldots,\gamma_{uK})A_0^{-1}. \tag{2}
\]
Because \(\Gamma_Y(x)\) is nonsingular, its rows span \(\mathbb R^K\).  Hence the diagonal matrices on the right of (2) span the entire diagonal algebra.  The common one-dimensional invariant subspaces of that algebra are precisely the coordinate axes, so the common one-dimensional invariant subspaces of the observable family in (2) are
\(\operatorname{span}(A_{0,\cdot j})\), \(j=1,\ldots,K\).  This recovers the columns of \(A_0\) up to separate nonzero scales and one common permutation.  Each column is a probability vector and satisfies \(\mathbf 1_K^{\mathsf T}A_{0,\cdot j}=1\), which fixes its scale (including its sign).

Set \(H=A_0^{-1}S\).  Since the columns of \(Z_0\) are probability vectors,
\[
H=\Pi Z_0^{\mathsf T},\qquad
H\mathbf 1_K=\Pi Z_0^{\mathsf T}\mathbf 1_K=\pi(x),\qquad
Z_0^{\mathsf T}=\Pi^{-1}H. \tag{3}
\]
Thus \(\pi(x)\) and \(Z_0\) are recovered with the same permutation.  Returning to (1) then gives
\[
\operatorname{diag}(\gamma_{u1},\ldots,\gamma_{uK})
=\Pi^{-1}A_0^{-1}T_u(Z_0^{\mathsf T})^{-1}, \tag{4}
\]
so \(\Gamma_Y(x)\) is recovered as well.

This recovery is measurable.  Indeed, choose from a fixed countable dense subset of \(\mathbb R^K\) the first vector \(r\) for which \(\sum_u r_uT_uS^{-1}\) has distinct eigenvalues.  Such an \(r\) exists because the columns of \(\Gamma_Y(x)\) are distinct, and the event that a candidate has distinct eigenvalues is Borel.  Eigenlines, inversion on the nonsingular matrices, and column normalization have Borel versions on this set.  The modal-label margin assigns each recovered column to the unique row where its largest entry occurs; the positive gap makes this permutation unique and measurable.  Hence the preceding display is an observed-law recovery map, not a definition of the latent factors.

Let \(p_A(y\mid x)\) be the observable \(K\)-vector whose \(a\)-th entry is the conditional joint density of \((Y,A=a)\) given \(X=x\).  Domination and conditional independence give, for \(\nu\)-almost every \(y\),
\[
p_A(y\mid x)=A_0\Pi f(y\mid x),\qquad
f(y\mid x)=\Pi^{-1}A_0^{-1}p_A(y\mid x). \tag{5}
\]
Every division and inversion in (2)--(5) is justified by the four stated positivity and rank conditions.  Equations (1)--(5), followed by modal alignment, therefore identify all objects asserted in the theorem.

It remains to relate the observed-law functional to the causal target.  Equation (5) implies
\(m_j(x)=E(Y\mid J=j,X=x)\).  Consistency and latent ignorability give, on the positive-propensity stratum,
\[
E\{Y(j)\mid J=j,X=x\}=E(Y\mid J=j,X=x)=m_j(x),
\qquad
E\{Y(j)\mid J=j,X=x\}=E\{Y(j)\mid X=x\}. \tag{6}
\]
The outcome is bounded, so these conditional means are integrable.  Taking \(P_X\)-expectations in (6) proves
\[
E\{Y(j)\}=E_P\{m_j(X)\}=\psi_j,
\]
for every \(j\in\mathcal J\), as required.
\end{proof}
\par\smallskip\noindent \textit{Justification.} This isolates point identification from the later efficiency claim and fixes the causal meaning of each recovered class. \textit{Gap.} Allman, Matias, and Rhodes identify factors up to labels, while Zhou and Tchetgen Tchetgen give the binary hidden-treatment version; the conditional measurable \(K\)-class formulation must be verified on the stated quantitative stratum. \textit{Consumer.} It makes the six maternal-smoking exposure levels interpretable as ordered intervention states rather than anonymous mixture classes.

\begin{proposition}[prop:2-product-bridge]\label{prop:2-product-bridge}
For every \(P\in\mathcal M_{K,\varepsilon}\), \(\phi_j=U_j-\psi_j\) is an observed-data influence function for \(\psi_j\); it is the static, zero-discount specialization of the categorical product-bridge construction in Appendix A.8 of Bian, Zhou, and Cui.
\end{proposition}
\begin{proof}
Fix \(j\in\mathcal J\).  The inverse-feature definitions, the conditional measurement laws, and Assumptions \(\mathrm{ass{:}surrogate\mbox{-}rank}\) and \(\mathrm{ass{:}proxy\mbox{-}rank}\) give
\[
\begin{aligned}
 E_P\{a_j(A,X)\mid J=\ell,X\}
 &=e_j^{\mathsf T}M_A(X)^{-1}M_A(X)e_\ell
 =\mathbf 1\{j=\ell\},\\
 E_P\{z_j(Z,X)\mid J=\ell,X\}
 &=e_j^{\mathsf T}M_Z(X)^{-1}M_Z(X)e_\ell
 =\mathbf 1\{j=\ell\}. \tag{1}
\end{aligned}
\]
By Assumption \(\mathrm{ass{:}three\mbox{-}view\mbox{-}independence}\), conditional on \((J=\ell,X)\) the three factors \(a_j(A,X)\), \(z_j(Z,X)\), and \(Y-m_j(X)\) are independent.  Hence
\[
 E_P(U_j\mid J=\ell,X)
 =m_j(X)+\frac{\mathbf 1\{j=\ell\}}{\pi_j(X)}
 E_P\{Y-m_j(X)\mid J=j,X\}=m_j(X). \tag{2}
\]
The last conditional expectation is zero by the definition of \(m_j\).  Thus \(E_P(U_j)=E_P\{m_j(X)\}=\psi_j\), and \(\phi_j=U_j-\psi_j\) is centered.

Consider an observed DQM path and let \(s\) be its observed score.  Theorem \(\mathrm{thm{:}differential\mbox{-}score\mbox{-}lifting}\) supplies factorwise centered scores \(s_X,s_\pi,s_Y,s_A,s_Z\) such that
\[
s=E_P(s_X+s_\pi+s_Y+s_A+s_Z\mid O).
\]
Repeated conditioning, (1), and factorwise centering give
\[
\begin{aligned}
 E_P(\phi_js_X)&=E_P[\{m_j(X)-\psi_j\}s_X(X)],\\
 E_P(\phi_js_\pi)&=0,\\
 E_P(\phi_js_Y)&=E_P\!\left[
 \frac{\mathbf 1\{J=j\}}{\pi_j(X)}
 \{Y-m_j(X)\}s_Y(Y,J,X)\right],\\
 E_P(\phi_js_A)&=E_P(\phi_js_Z)=0. \tag{3}
\end{aligned}
\]
For the propensity line, use (2) and \(E_P(s_\pi\mid X)=0\).  For the \(A\)-score line, \(m_j(X)-\psi_j\) vanishes against \(E_P(s_A\mid J,X)=0\).  In the residual term, \(\ell=j\) gives \(E_P(Y-m_j\mid J=j,X)=0\), whereas \(\ell\ne j\) gives \(E_P(z_j\mid J=\ell,X)=0\).  The \(Z\)-score line is symmetric.  For the outcome line, both bridges average to \(\mathbf 1\{J=j\}\) by (1).

The first line of (3) is the derivative of the outer expectation over \(P_X\).  Assumption \(\mathrm{ass{:}latent\mbox{-}positivity}\) permits division by \(\pi_j(X)\), and
\[
 E_P\!\left[
 \frac{\mathbf 1\{J=j\}}{\pi_j(X)}
 \{Y-m_j(X)\}s_Y\,\middle|\,X\right]
 =E_P[\{Y-m_j(X)\}s_Y\mid J=j,X], \tag{4}
\]
which is the derivative of the class-\(j\) conditional mean.  The target is unchanged by propensity and measurement-factor perturbations, agreeing with the zero lines of (3).  Since \(\phi_j\) is \(O\)-measurable, conditional-expectation duality yields
\[
 E_P(\phi_js)
 =E_P\{\phi_j(s_X+s_\pi+s_Y+s_A+s_Z)\}. \tag{5}
\]
Equations (3)--(5) equal the pathwise derivative of \(\psi_j\) along every observed DQM path, proving the influence-function equation.

Assumption \(\mathrm{ass{:}bounded\mbox{-}outcome}\) gives \(|m_j(X)|\leq L\).  Latent positivity gives \(\pi_j^{-1}\leq\varepsilon^{-1}\), and the two measurement-rank assumptions give
\(\|M_A^{-1}\|_{\mathrm{op}}\vee\|M_Z^{-1}\|_{\mathrm{op}}\leq\varepsilon^{-1}\).
The category vectors have unit norm, so \(\phi_j\) is bounded by a constant depending only on the regular stratum.  Therefore \(\phi_j\in L^2_0(P)\).

Finally, the cited leaf \(\mathrm{lem{:}bian\mbox{-}categorical\mbox{-}product\mbox{-}bridge}\) supplies the categorical inverse-feature and pair-product construction.  In its categorical-action formula, take the state to be \(X\), set the discount to zero, use the same initial and observed state law \(P_X\), let the target policy put unit mass on class \(j\), take behavior propensity \(\pi_j(X)\), reward \(Y\), and inverse features \(a_j\) and \(z_j\).  Its displayed score then equals \(U_j-\psi_j\).  This proves both the stated specialization and its validity here.
\end{proof}
\par\smallskip\noindent \textit{Justification.} The finite projection needs a valid starting gradient whose relationship to published categorical bridges is exact. \textit{Gap.} Bian, Zhou, and Cui provide validity but not canonicality for the present static continuous-outcome target. \textit{Consumer.} It lets analysts retain the multiply robust bridge architecture while replacing its avoidable variance by the canonical correction.

\begin{lemma}[lem:2a-constant-x-feature-inverse]\label{lem:2a-constant-x-feature-inverse}
For every \(P\in\mathcal M_{K,\varepsilon}\) under which \(X\) is degenerate, the derivative of the augmented normalized moment map \((S,T_1,\ldots,T_K)\), with \(S=P(A,Z)\) and \(T_u=E\{h_{0,u}(Y)e_Ae_Z^{\mathsf T}\}\), has a right inverse on its observed finite-feature score image whose operator norm is bounded by \(C_{\mathrm{reg}}\).  Applying the corrected \(\mathfrak I_0\) recovers centered derivatives, and hence centered scores, for \(\pi\), \(M_A\), \(M_Z\), and \(\Gamma_Y\) under the modal labeling.
\end{lemma}
\begin{proof}
Because \(X\) is degenerate, suppress it.  Put \(A_0=M_A\), \(Z_0=M_Z\), \(\Pi=\operatorname{diag}(\pi)\), and let \(\gamma_u^{\mathsf T}\) be row \(u\) of \(\Gamma_Y\).  Assumption \(\mathrm{ass{:}three\mbox{-}view\mbox{-}independence}\) gives
\[
 S=A_0\Pi Z_0^{\mathsf T},\qquad
 T_u=A_0\Pi\operatorname{diag}(\gamma_u)Z_0^{\mathsf T}. \tag{6}
\]
Consider a derivative \((\dot\pi,\dot A_0,\dot Z_0,\dot\Gamma)\) in normalized tangent coordinates:
\[
 \mathbf 1_K^{\mathsf T}\dot\pi=0,\qquad
 \mathbf 1_K^{\mathsf T}\dot A_0=0,\qquad
 \mathbf 1_K^{\mathsf T}\dot Z_0=0. \tag{7}
\]
We first prove injectivity of the derivative of (6).  Suppose \(\dot S=0\) and \(\dot T_u=0\) for every \(u\).  Define
\[
 E_A=A_0^{-1}\dot A_0,\qquad E_Z=Z_0^{-1}\dot Z_0,
 \qquad p_j=\dot\pi_j/\pi_j. \tag{8}
\]
The inverses and the divisions are valid by the surrogate-rank, proxy-rank, and latent-positivity assumptions.  Left-multiplying \(\dot S=0\) by \(A_0^{-1}\), right-multiplying by \(Z_0^{-\mathsf T}\), and taking entry \((r,s)\), \(r\ne s\), gives
\[
 (E_A)_{rs}\pi_s+\pi_r(E_Z)_{sr}=0. \tag{9}
\]
Applying the same transformations to \(\dot T_u=0\) and subtracting \(\gamma_{ur}\) times (9) gives
\[
 (E_A)_{rs}\pi_s(\gamma_{us}-\gamma_{ur})=0,
 \qquad u=1,\ldots,K. \tag{10}
\]
The outcome-feature rank assumption implies
\[
 \|\gamma_s-\gamma_r\|_2
 =\|\Gamma_Y(e_s-e_r)\|_2
 \geq\varepsilon\|e_s-e_r\|_2=\sqrt 2\,\varepsilon. \tag{11}
\]
Together with \(\pi_s\geq\varepsilon\), (10)--(11) imply \((E_A)_{rs}=0\); (9) then implies \((E_Z)_{sr}=0\).  Thus \(E_A\) and \(E_Z\) are diagonal.  Since \(\mathbf 1_K^{\mathsf T}A_0=\mathbf 1_K^{\mathsf T}\), (7)--(8) give \(\mathbf 1_K^{\mathsf T}E_A=0\), so the diagonal matrix \(E_A\) is zero.  Similarly \(E_Z=0\).  The diagonal entries of the transformed equation \(\dot S=0\) now give \(\dot\pi=0\), and those of every transformed equation \(\dot T_u=0\) give \(\dot\gamma_{ur}=0\) for all \(u,r\).  The derivative is therefore injective on the normalized tangent space.

The admissible finite tuples lie in a compact set.  Indeed, the stochastic columns and \(\pi\) lie in finite products of compact simplices; the least-singular-value restrictions and propensity lower bound define closed subsets; and \(|\Gamma_{Y,uj}|\leq\|h_0\|_\infty\).  The derivative matrix in (6) is continuous in this tuple.  Its least singular value on the normalized tangent space is positive pointwise by the preceding injectivity argument.  Were its infimum zero, compactness would give a convergent minimizing subsequence and a unit null vector at its limit, contradicting injectivity.  Hence the infimum is a positive constant depending only on \(K,\varepsilon\), and \(\|h_0\|_\infty\).  The Moore--Penrose right inverse on the observed image is therefore bounded by \(C_{\mathrm{reg}}\), after enlarging that declared constant if necessary.

On a fixed-rank finite-dimensional set, the Moore--Penrose inverse is Borel measurable: on each chart where a fixed maximal minor is nonzero, the inverse is a ratio of determinants, and finitely many charts cover the set.  Theorem \(\mathrm{thm{:}1\mbox{-}identification}\) supplies the identified common permutation, while the modal gap makes its ordered choice locally constant.  Applying this Borel inverse to the differentiated observed moments therefore recovers the normalized derivatives measurably under the ordered labeling.  At a positive categorical atom, its factor score is the derivative divided by the atom probability.  At a zero atom, nonnegativity of a two-sided probability path forces its derivative to be zero, so assigning score zero is valid.  Equation (7) gives factorwise centering.  This proves the asserted bounded right inverse on the observed finite-feature score image.
\end{proof}
\par\smallskip\noindent \textit{Justification.} This finite-dimensional derivative inversion is the mandatory first gate before any measurable conditional-\(X\) or full-density lifting claim. \textit{Gap.} Allman, Matias, and Rhodes prove level uniqueness of discrete factors, while Hu and Schennach prove full-law recovery through operators; neither result states this quantitative right inverse for the finite normalized tensor derivative. \textit{Consumer.} It provides the constant-covariate kill test and the local inverse needed to begin the full tangent-completeness proof.

\begin{theorem}[thm:2-tangent-completeness]\label{thm:2-tangent-completeness}
At every two-sided pathwise-interior \(P\in\mathcal M_{K,\varepsilon}\), meaning that every bounded factorwise centered additive latent score generates for all sufficiently small positive and negative path parameters a product-factor path remaining in \(\mathcal M_{K,\varepsilon}\), the full observed-data tangent space satisfies \(\mathcal T=\overline{\mathcal S_{\mathrm{lat}}}^{\,L^2(P)}\).  At such a \(P\),
\[
\mathcal T^{\perp}=\{k\in L^2_0(P):E(k\mid Y,J,X)=E(k\mid A,J,X)=E(k\mid Z,J,X)=0\}.
\]
\end{theorem}
\begin{proof}
Theorem \(\mathrm{thm{:}differential\mbox{-}inverse\mbox{-}resolution}\) places every observed DQM score in \(\mathcal S_{\mathrm{lat}}\).  Taking \(L^2(P)\) closures and using the definition of \(\mathcal T\) gives
\[
 \mathcal T\subseteq\overline{\mathcal S_{\mathrm{lat}}}^{\,L^2(P)}. \tag{12}
\]
Conversely, at the stated two-sided pathwise-interior point, every bounded factorwise centered additive latent score generates, for all sufficiently small parameters of either sign, a product-factor DQM path in \(\mathcal M_{K,\varepsilon}\).  Differentiating the product likelihood and marginalizing \(J\) shows that the observed score is the conditional expectation of the additive latent score given \(O\), and hence belongs to \(\mathcal T\).  For an arbitrary square-integrable factor score, truncate it and subtract its conditional factor mean.  Conditional Jensen's inequality gives convergence to the original score in the corresponding factor \(L^2\) norm.  Conditional expectation is an \(L^2\) contraction, so the associated observed scores also converge.  Since \(\mathcal T\) is closed, this proves the reverse inclusion in (12), and therefore
\[
 \mathcal T=\overline{\mathcal S_{\mathrm{lat}}}^{\,L^2(P)}. \tag{13}
\]

It remains to calculate the orthocomplement.  Let \(k\in L^2_0(P)\).  Orthogonality to every centered outcome-factor score is equivalent to
\[
 E_P(k\mid Y,J,X)=E_P(k\mid J,X). \tag{14}
\]
For necessity, subtract the right side from the left; the difference is itself a centered outcome-factor function, and testing against that difference makes its squared norm zero.  Sufficiency follows by iterated expectation.  The identical argument with \(A\) and \(Z\) gives the analogous equalities.  Orthogonality to every centered propensity score is equivalent to
\(E_P(k\mid J,X)=E_P(k\mid X)\), and orthogonality to every centered \(X\)-marginal score is equivalent to \(E_P(k\mid X)=E_P(k)=0\).  Thus orthogonality to \(\mathcal S_{\mathrm{lat}}\) implies
\[
 E_P(k\mid Y,J,X)=E_P(k\mid A,J,X)=E_P(k\mid Z,J,X)=0. \tag{15}
\]
Conversely, the first equality in (15) implies by iteration that \(E_P(k\mid J,X)=E_P(k\mid X)=0\).  Each inner product with an \(X\)- or propensity-factor score is then zero.  The three displayed equalities make the inner products with the outcome and two measurement factor scores zero as well.  Hence \(k\perp\mathcal S_{\mathrm{lat}}\), and (13) yields \(k\perp\mathcal T\).  This proves the displayed characterization of \(\mathcal T^\perp\).
\end{proof}
\par\smallskip\noindent \textit{Justification.} Equality with the full observed DQM closure, not merely a convenient submodel score span, is the key efficiency theorem. \textit{Gap.} Continuous full-law recovery in Hu and Schennach and its triple-proxy causal use in Deaner are identification statements, not a measurable uniformly bounded inverse of the derivative of the factor map. The binary tangent calculation in Zhou and Tchetgen Tchetgen also does not establish that multilevel DQM lifting on the stated regular stratum. \textit{Consumer.} It turns the candidate variance calculation into the information bound for every regular estimator of the smoking-level contrasts.

\begin{lemma}[lem:3a-projector-algebra]\label{lem:3a-projector-algebra}
Under common support, latent positivity, and observed-cell positivity, \(W\) and \(D\) are invertible; the surrogate and proxy rank conditions imply that \(B\) has full column rank. Bayes' rule gives \(R=D^{-1}BW\), \(R^{\mathsf T}D=WB^{\mathsf T}\), and \(R^{\mathsf T}DR=WB^{\mathsf T}D^{-1}BW\). Hence \(N_B=I_{K^{2}}-D^{-1}B(B^{\mathsf T}D^{-1}B)^{+}B^{\mathsf T}=I_{K^{2}}-R(R^{\mathsf T}DR)^{+}R^{\mathsf T}D\), without any additional restriction on \(w\), and \(N_B^2=N_B\), \(N_B^{\mathsf T}D=DN_B\), \(B^{\mathsf T}N_B=0\), \(\operatorname{range}(N_B)=\ker(B^{\mathsf T})\), and \(\ker(N_B)=\operatorname{range}(D^{-1}B)=\operatorname{range}(R)\). In the exact symbolic configuration \(\mathfrak A_3\) with \(W=\operatorname{diag}(w_1,w_2,w_3)\) and positive nonconstant \(w_j\), all six identities hold by matrix multiplication for every full-column-rank \(B\) with \(Bw>0\).
\end{lemma}
\begin{proof}
Latent positivity and common support give \(w_j(y,x)>0\), hence \(W\) is invertible; cell positivity makes \(D\) invertible.  To prove that \(B\) has full column rank, suppose \(Bc=0\).  Reshaping this vector into a \(K\times K\) matrix gives
\[
M_A\operatorname{diag}(c)M_Z^{\mathsf T}=0.
\]
Multiplication by the two inverses, licensed by the surrogate and proxy rank assumptions, gives \(\operatorname{diag}(c)=0\), hence \(c=0\).

Bayes' rule and the conditional product factorization give, for \(r=(a,z)\),
\[
R_{rj}=\frac{w_jB_{rj}}{d_r},
\qquad R=D^{-1}BW,
\qquad R^{\mathsf T}D=WB^{\mathsf T},
\qquad R^{\mathsf T}DR=WB^{\mathsf T}D^{-1}BW. \tag{18}
\]
Let \(H=B^{\mathsf T}D^{-1}B\).  The established ranks make \(H\) and \(R^{\mathsf T}DR=WHW\) positive definite, so every displayed pseudoinverse is its ordinary inverse.  Therefore
\[
R(R^{\mathsf T}DR)^{-1}R^{\mathsf T}D
=D^{-1}BW(W^{-1}H^{-1}W^{-1})WB^{\mathsf T}
=D^{-1}BH^{-1}B^{\mathsf T}. \tag{19}
\]
This proves the equality of the two formulas for \(N_B\), without any equality restriction on the entries of \(w\).

Writing \(P_B=D^{-1}BH^{-1}B^{\mathsf T}\), direct multiplication gives \(P_B^2=P_B\), \(P_B^{\mathsf T}D=DP_B\), and \(B^{\mathsf T}(I-P_B)=0\).  Hence
\[
N_B^2=N_B,\qquad N_B^{\mathsf T}D=DN_B,\qquad B^{\mathsf T}N_B=0. \tag{20}
\]
If \(B^{\mathsf T}v=0\), then \(N_Bv=v\), proving \(\operatorname{range}N_B=\ker B^{\mathsf T}\).  Also \(N_BD^{-1}B=0\), and both \(\operatorname{range}(D^{-1}B)\) and \(\ker N_B\) have dimension \(K\); thus they are equal.  Equation (18) and invertibility of \(W\) show that this space also equals \(\operatorname{range}R\).  In \(\mathfrak A_3\), the same proof uses ordinary inverses because its hypotheses give positive \(W,D\) and full-column-rank \(B\); (18)--(20) are the requested six exact symbolic identities for arbitrary positive nonconstant \(w\).
\end{proof}
\par\smallskip\noindent \textit{Justification.} This algebra fixes the first outcome-score elimination and independently stress-tests it under nonuniform posterior weights. \textit{Gap.} The posterior-coordinate formula is valid only after the Bayes reparameterization; the required target is \(\ker(B^{\mathsf T})\), not the distinct cell-weighted null space \(\ker(B^{\mathsf T}D)\). \textit{Consumer.} It supplies an exact symbolic unit test for every implementation of the nine-cell \(K=3\) operator.

\begin{theorem}[thm:3-finite-projector]\label{thm:3-finite-projector}
For every two-sided pathwise-interior \(P\in\mathcal M_{K,\varepsilon}\), the outcome-score restrictions \(E_P(k\mid Y,J=j,X)=0\) are equivalent, cellwise, to \(B^{\mathsf T}k=0\).  The operator \(N=N_B=I_{K^2}-D^{-1}B(B^{\mathsf T}D^{-1}B)^+B^{\mathsf T}\) is their pointwise \(D\)-orthogonal projector, \(b_{f_0}(X)\in\operatorname{range}\{G(X)\}\), and \(Qf_0=N_Bf_0-h^{\mathsf T}G(X)^+b_{f_0}(X)\) is the \(L^2(P)\)-orthogonal projection onto \(\mathcal T^\perp\) for every \(f_0\in L^2(P)\).
\end{theorem}
\begin{proof}
Fix \((y,x)\) outside the null sets in the assumptions and represent an observed-cell function by \(k=(k_{az})\in\mathbb R^{K^2}\).  Conditional three-view independence gives
\[
 E_P\{k(O)\mid Y=y,J=j,X=x\}
 =\sum_{a,z}M_A(x)_{aj}M_Z(x)_{zj}k_{az}
 =B_{\cdot j}(x)^{\mathsf T}k. \tag{16}
\]
Thus the outcome-score restrictions in Theorem \(\mathrm{thm{:}2\mbox{-}tangent\mbox{-}completeness}\) are exactly \(B^{\mathsf T}k=0\).  Lemma \(\mathrm{lem{:}3a\mbox{-}projector\mbox{-}algebra}\) proves that \(N_B\) is idempotent, \(D\)-self-adjoint, and has range \(\ker(B^{\mathsf T})\); hence it is the pointwise \(D\)-orthogonal projector onto the outcome restrictions.  Moreover \(R\mathbf 1_K=\mathbf 1_{K^2}\), because every posterior row sums to one.  Lemma 3A then gives \(N_B\mathbf 1_{K^2}=0\), so \(N_Bf_0\) has conditional mean zero given \((Y,X)\).

For a coordinate \(r=(j,a)\), Bayes' identity and iterated expectation give
\[
\begin{aligned}
 E_P\{kR_{\cdot j}\mathbf 1(A=a)\mid X\}
 &=E_P\{k\mathbf 1(J=j)\mathbf 1(A=a)\mid X\}\\
 &=\pi_j(X)E_P\{k\mathbf 1(A=a)\mid J=j,X\}. \tag{17}
\end{aligned}
\]
The analogous identity holds for \(r=(j,z)\).  Since \(\pi_j\geq\varepsilon\), the remaining measurement restrictions in Theorem 2 are exactly \(E_P(kc_r\mid X)=0\) for all \(r\), with zero-probability measurement atoms imposing no restriction.  If \(k=N_Bk\), conditional \(D\)-self-adjointness yields
\[
 E_P(kc_r\mid X)=E_P\{k(N_Bc_r)\mid X\}=E_P(kh_r\mid X). \tag{18}
\]

If \(v\in\ker G(X)\), then
\[
 0=v^{\mathsf T}G(X)v=E_P\{(v^{\mathsf T}h)^2\mid X\}, \tag{19}
\]
so \(v^{\mathsf T}h=0\) conditionally and \(v^{\mathsf T}b_{f_0}=0\).  Hence
\(b_{f_0}\perp\ker G=\operatorname{range}(G)^\perp\), and therefore \(b_{f_0}\in\operatorname{range}(G)\).  It follows that
\[
 E_P\{h(Qf_0)\mid X\}
 =b_{f_0}-GG^+b_{f_0}=0. \tag{20}
\]
Every coordinate of \(h=N_Bc\) is in \(\operatorname{range}(N_B)\), so \(Qf_0\) remains in that range.  Equations (16)--(20) and Theorem 2 imply \(Qf_0\in\mathcal T^\perp\).

For the projection property, take \(k\in\mathcal T^\perp\).  Pointwise \(D\)-orthogonality gives
\[
 E_P[((I-N_B)f_0)k\mid Y,X]=0. \tag{21}
\]
Equation (18) gives \(E_P(hk\mid X)=0\), whence
\[
 E_P\{h^{\mathsf T}G^+b_{f_0}\,k\}
 =E_P[(G^+b_{f_0})^{\mathsf T}E_P(hk\mid X)]=0. \tag{22}
\]
Thus \(E_P\{(f_0-Qf_0)k\}=0\).  The space characterized by the three conditional-zero restrictions is an intersection of kernels of continuous conditional-expectation maps and is closed.  Therefore \(Q\) is the \(L^2(P)\)-orthogonal projection onto \(\mathcal T^\perp\).

Assumption \(\mathrm{ass{:}cell\mbox{-}positivity}\) bounds conversion between the conditional \(D\)-norm and Euclidean cell norm.  The posterior and projection Gram gaps bound the displayed inverses.  Constant rank makes the pseudoinverse Borel measurable by the finite-minor chart argument used in Lemma 2A.  Conditional Cauchy--Schwarz then shows that every displayed term lies in \(L^2(P)\), completing the proof.
\end{proof}
\par\smallskip\noindent \textit{Justification.} The theorem eliminates the infinite continuous-outcome nuisance pointwise and leaves only one finite conditional Gram solve. \textit{Gap.} Zhou and Tchetgen Tchetgen use an equal-width outcome partition and nuisance-dependent conditional-probability functions for the binary continuous-outcome projection; no located source gives this exact multilevel \(K^{2}\)-cell plus \(2K^{2}\)-Gram representation. \textit{Consumer.} It makes exact efficiency computation feasible without selecting an outcome sieve in the continuous birth-weight application.

\begin{theorem}[thm:4-canonical-gradient]\label{thm:4-canonical-gradient}
At every two-sided pathwise-interior \(P\in\mathcal M_{K,\varepsilon}\) and every fixed \(K\geq3\), the vector \(\phi^{\mathrm{eff}}=\phi-Q\phi\), with first elimination \(N_B=I_{K^2}-D^{-1}B(B^{\mathsf T}D^{-1}B)^+B^{\mathsf T}\), is the unique canonical gradient for \(\psi\), and \(V_{\mathrm{eff}}\) is the semiparametric efficiency bound.  The operator uses fixed \(K^2\)-cell algebra and one Gram solve of dimension at most \(2K^2\), with no outcome-partition resolution or outcome-basis approximation index.
\end{theorem}
\begin{proof}
Proposition \(\mathrm{prop{:}2\mbox{-}product\mbox{-}bridge}\) gives an observed influence-function vector \(\phi\).  Theorem \(\mathrm{thm{:}3\mbox{-}finite\mbox{-}projector}\) says that \(Q\phi\) is its orthogonal projection onto \(\mathcal T^\perp\).  Hence
\[
 \phi^{\mathrm{eff}}=(I-Q)\phi\in\mathcal T. \tag{23}
\]
For every score \(s\in\mathcal T\), orthogonality gives
\[
 E_P(\phi^{\mathrm{eff}}s)
 =E_P(\phi s)-E_P\{(Q\phi)s\}
 =\dot\psi(s). \tag{24}
\]
Thus \(\phi^{\mathrm{eff}}\) remains an influence function.  If \(\widetilde\phi\in\mathcal T\) is another influence function, the influence equations imply
\(\widetilde\phi-\phi^{\mathrm{eff}}\in\mathcal T^\perp\).  The difference also belongs to \(\mathcal T\), and \(\mathcal T\cap\mathcal T^\perp=\{0\}\), so uniqueness follows.

Every other influence function can be written \(\phi^{\mathrm{eff}}+r\) with \(r\in\mathcal T^\perp\).  Orthogonality yields, for every \(c\in\mathbb R^K\),
\[
 \operatorname{Var}_P\{c^{\mathsf T}(\phi^{\mathrm{eff}}+r)\}
 =\operatorname{Var}_P(c^{\mathsf T}\phi^{\mathrm{eff}})
  +E_P\{(c^{\mathsf T}r)^2\}. \tag{25}
\]
Thus \(V_{\mathrm{eff}}\) is the semiparametric efficiency bound.

The first projection acts on exactly the \(K^2\) observed \((A,Z)\)-cells at each \((Y,X)\).  The second uses the Gram matrix of the displayed \(2K^2\) residualized measurement coordinates and therefore requires one solve of dimension at most \(2K^2\), possibly smaller after rank reduction.  Dependence on continuous \(Y\) is pointwise or through conditional expectation; no outcome partition or outcome-basis approximation is introduced.  This proves the computational-dimension assertion.
\end{proof}
\par\smallskip\noindent \textit{Justification.} This is the theorem-level kernel: an exact continuous-outcome efficiency bound with fixed finite linear algebra. \textit{Gap.} The closest result defines the binary continuous canonical gradient abstractly and approximates its projection; the exact multilevel projector and tangent-completeness proof are absent. \textit{Consumer.} It supplies the variance-minimizing regular score for all prespecified smoking-intensity contrasts.

\begin{theorem}[thm:4b-partition-comparison]\label{thm:4b-partition-comparison}
At every two-sided pathwise-interior law in the binary restriction \(J^{\mathrm{bin}}\in\{0,1\}\), let \(\mathcal I^{\ell}\), \(\mathcal Q^{\ell}\), \(\mathcal P_1^{\ell}\), the nuisance-dependent conditional-probability functions \(p_1^{\ell},\ldots,p_{\ell-2}^{\ell}\), and \(V^{\ell}=(B_0p_1^{\ell},B_1p_1^{\ell},\ldots,B_0p_{\ell-2}^{\ell},B_1p_{\ell-2}^{\ell})^{\mathsf T}\) be exactly the outcome-partition construction of Zhou and Tchetgen Tchetgen Supplement S8.2. With \(Q\) built from the corrected \(D\)-orthogonal first elimination \(N_B\), every well-defined \(\operatorname{span}(V^{\ell})\) is contained in \(\operatorname{range}(Q)\), and \(V_{\ell}^{\mathrm{ZTT}}-V_{\mathrm{eff}}=E_P[\{(Q-Q_{\ell}^{\mathrm{ZTT}})\phi\}\{(Q-Q_{\ell}^{\mathrm{ZTT}})\phi\}^{\mathsf T}]\succeq0\). Here \(V_{\mathrm{eff}}\) is the model efficiency bound only at such interior laws. The gap is strict in a contrast direction \(c\) exactly when \(c^{\mathsf T}(Q-Q_{\ell}^{\mathrm{ZTT}})\phi\neq0\). Along a nested equal-width refinement subsequence, the Supplement S8.2 density of \(\mathcal P_1^{\ell}\) in the binary conditional-null space implies \(Q_{\ell}^{\mathrm{ZTT}}\phi\to Q\phi\) in \(L^{2}(P)\) and \(V_{\ell}^{\mathrm{ZTT}}\to V_{\mathrm{eff}}\). This comparison is a binary restriction only; the exact fixed-cell theorem is asserted for \(K\geq3\).
\end{theorem}
\begin{proof}
Fix a two-sided pathwise-interior law in the binary restriction and a value of \(\ell\) for which every denominator in the displayed basis and the matrix \(E_P(V^\ell V^{\ell,\mathsf T})\) are invertible.  All Hilbert-space statements below refer to \(L^2(P)\), and every operator applied to a vector is applied coordinatewise.

Let
\[
\mathcal H=\operatorname{range}(Q),
\qquad
\mathcal H_\ell=\operatorname{span}\{V_1^\ell,\ldots,V_{2(\ell-2)}^\ell\}.
\]
Proposition 10 identifies \(\mathcal H\) in the binary restriction as the closed span of functions
\[
B_0h_0(Y,X)+B_1h_1(Y,X),
\qquad
E_P\{h_r(Y,X)\mid J^{\mathrm{bin}},X\}=0,quad r\in\{0,1\},                                      \tag{1}
\]
and identifies \(Q\) as the orthogonal projection onto this observed tangent-space orthocomplement.  By the definition of \(\mathcal P_1^\ell\) and the Supplement S8.2 basis construction, each \(p_t^\ell\) belongs to \(\mathcal P_1^\ell\), and therefore
\[
E_P\{p_t^\ell(Y,X)\mid J^{\mathrm{bin}},X\}=0.                                            \tag{2}
\]
The coordinates \(B_0p_t^\ell\) and \(B_1p_t^\ell\) are consequently instances of (1), obtained by setting, respectively, \((h_0,h_1)=(p_t^\ell,0)\) and \((h_0,h_1)=(0,p_t^\ell)\).  Thus
\[
\mathcal H_\ell\subseteq\mathcal H.                                                       \tag{3}
\]
This proves the asserted finite-\(\ell\) range inclusion.  The corrected formula for \(N_B\) is legitimate here by Lemma 3A; in particular, the inclusion uses the \(D\)-orthogonal elimination whose range is \(\ker(B^{\mathsf T})\), rather than a different cell-weighted null space.

Put \(\Gamma_\ell=E_P(V^\ell V^{\ell,\mathsf T})\).  For scalar \(f\in L^2(P)\), the comparator definition gives
\[
Q_\ell^{\mathrm{ZTT}}f
=E_P(fV^{\ell,\mathsf T})\Gamma_\ell^{-1}V^\ell.                                           \tag{4}
\]
For any \(u=\alpha^{\mathsf T}V^\ell\in\mathcal H_\ell\), (4) yields
\[
E_P\!\left[\{f-Q_\ell^{\mathrm{ZTT}}f\}u\right]
=E_P(fV^{\ell,\mathsf T})\alpha
-E_P(fV^{\ell,\mathsf T})\Gamma_\ell^{-1}\Gamma_\ell\alpha
=0.                                                                                        \tag{5}
\]
Moreover, \(Q_\ell^{\mathrm{ZTT}}f\in\mathcal H_\ell\).  Hence (4)--(5) prove directly that \(Q_\ell^{\mathrm{ZTT}}\) is the orthogonal projection onto \(\mathcal H_\ell\).  From (3) and the corresponding projection property of \(Q\),
\[
Q_\ell^{\mathrm{ZTT}}Q=QQ_\ell^{\mathrm{ZTT}}=Q_\ell^{\mathrm{ZTT}},
\qquad
(Q-Q_\ell^{\mathrm{ZTT}})^2=Q-Q_\ell^{\mathrm{ZTT}},                                      \tag{6}
\]
and \(Q-Q_\ell^{\mathrm{ZTT}}\) is the orthogonal projection onto \(\mathcal H\cap\mathcal H_\ell^\perp\).

Because \(Q\) projects onto the observed tangent-space orthocomplement and \(\phi\) is an influence-function vector, \(\phi-Q\phi\) is the canonical gradient at the stipulated two-sided interior law.  Indeed, subtracting \(Q\phi\), which is orthogonal to every tangent score, leaves all pathwise derivative equations unchanged, while \((I-Q)\phi\) belongs to the tangent-space closure.  If another element of that closure represented the same derivatives, its difference from \((I-Q)\phi\) would belong both to the closure and to its orthogonal complement, and hence would be zero.  This also explains why the efficiency interpretation is asserted here only at the stated pathwise-interior laws.  Therefore
\[
V_{\mathrm{eff}}=E_P\!\left[\{(I-Q)\phi\}\{(I-Q)\phi\}^{\mathsf T}\right].               \tag{7}
\]

Now decompose
\[
\phi-Q_\ell^{\mathrm{ZTT}}\phi
=(I-Q)\phi+(Q-Q_\ell^{\mathrm{ZTT}})\phi.                                                   \tag{8}
\]
The first summand in (8) lies coordinatewise in \(\mathcal H^\perp\), whereas the second lies coordinatewise in \(\mathcal H\), by (6).  Every cross-moment therefore vanishes.  Taking outer products in (8), then expectations, and using (7) gives the exact identity
\[
V_\ell^{\mathrm{ZTT}}-V_{\mathrm{eff}}
=E_P\!\left[\{(Q-Q_\ell^{\mathrm{ZTT}})\phi\}
             \{(Q-Q_\ell^{\mathrm{ZTT}})\phi\}^{\mathsf T}\right]\succeq0.              \tag{9}
\]
For every fixed contrast vector \(c\),
\[
c^{\mathsf T}(V_\ell^{\mathrm{ZTT}}-V_{\mathrm{eff}})c
=E_P\!\left[\{c^{\mathsf T}(Q-Q_\ell^{\mathrm{ZTT}})\phi\}^2\right]
=\left\|c^{\mathsf T}(Q-Q_\ell^{\mathrm{ZTT}})\phi\right\|_{L^2(P)}^2.                  \tag{10}
\]
Thus the variance gap is strict in direction \(c\) if and only if \(c^{\mathsf T}(Q-Q_\ell^{\mathrm{ZTT}})\phi\) is nonzero in \(L^2(P)\), equivalently nonzero on a set of positive \(P\)-probability.

Finally, take a nested equal-width refinement subsequence, indexed again by \(\ell\), on which the displayed denominators and Gram inverses exist.  Then \(\mathcal H_\ell\) is increasing.  Lemma \(\mathrm{ztt\mbox{-}partition\mbox{-}density}\), the cited Supplement S8.2 density statement, together with Proposition 10, gives
\[
\overline{\bigcup_\ell\mathcal H_\ell}=\mathcal H.                                        \tag{11}
\]
For a scalar coordinate \(f\) of \(\phi\), write \(u=Qf\in\mathcal H\).  Since \(Q_\ell^{\mathrm{ZTT}}(I-Q)f=0\) by (3),
\[
Q_\ell^{\mathrm{ZTT}}f=Q_\ell^{\mathrm{ZTT}}u.                                            \tag{12}
\]
Given \(\delta>0\), (11) supplies an index \(r\) and \(v\in\mathcal H_r\) such that \(\|u-v\|_2<\delta\).  For every later index \(\ell\), nestedness gives \(v\in\mathcal H_\ell\); the best-approximation property established in (5) then yields
\[
\|Qf-Q_\ell^{\mathrm{ZTT}}f\|_2
=\|u-Q_\ell^{\mathrm{ZTT}}u\|_2
\leq\|u-v\|_2<\delta.                                                                     \tag{13}
\]
Applying (13) to each of the finitely many coordinates proves
\[
Q_\ell^{\mathrm{ZTT}}\phi\longrightarrow Q\phi
\quad\text{in }L^2(P).                                                                      \tag{14}
\]
If \(\Delta_\ell=(Q-Q_\ell^{\mathrm{ZTT}})\phi\), then (14) says \(\|\Delta_{\ell,j}\|_2\to0\) for every coordinate.  Cauchy--Schwarz gives
\[
\left|E_P(\Delta_{\ell,j}\Delta_{\ell,k})\right|
\leq\|\Delta_{\ell,j}\|_2\,\|\Delta_{\ell,k}\|_2\longrightarrow0.                    \tag{15}
\]
Equations (9) and (15) imply \(V_\ell^{\mathrm{ZTT}}\to V_{\mathrm{eff}}\) entrywise, and hence in every matrix norm because the binary dimension is fixed.  Every construction used in (1)--(15) is within the binary restriction.  The limiting comparison therefore neither proves nor replaces the exact fixed-\(K^2\)-cell result asserted for fixed \(K\geq3\).
\end{proof}
\par\smallskip\noindent \textit{Justification.} The theorem compares the exact projection only with the comparator's actual partition-dependent, nuisance-dependent construction and separates that binary audit from the \(K\geq3\) contribution. \textit{Gap.} Supplement S8.2 proves density of its equal-width partition spaces and proposes the finite-\(\ell\) projection, but it does not reduce the multilevel continuous projection to fixed \(K^{2}\)-cell and \(2K^{2}\)-Gram algebra. \textit{Consumer.} It lets a binary-analysis user quantify the residual variance of the published discretization while preventing that binary approximation from being misrepresented as an arbitrary fixed-basis module.

\begin{theorem}[thm:5-strict-variance]\label{thm:5-strict-variance}
For every adjacent maternal-smoking intensity contrast \(c_j=e_{j+1}-e_j\in\mathcal C_{\mathrm{adj}}\), compute \(N_Bv_j\) with \(N_B=I_{K^{2}}-D^{-1}B(B^{\mathsf T}D^{-1}B)^{+}B^{\mathsf T}\). The recovered augmented conditional Gram matrix \(G_j^{\mathrm{aug}}(X)\) satisfies \(\operatorname{rank}\{G_j^{\mathrm{aug}}(X)\}\leq\operatorname{rank}\{G(X)\}+1\), and the observable condition \(P_X[\operatorname{rank}\{G_j^{\mathrm{aug}}(X)\}=\operatorname{rank}\{G(X)\}+1]>0\), equivalently \(P_X\{\Lambda_j(X)>0\}>0\), certifies \(c_j^{\mathsf T}Q\phi\neq0\) and strict variance reduction for that dose step. As a supporting consequence, \(c_j^{\mathsf T}(V_{\mathrm{raw}}-V_{\mathrm{eff}})c_j=E_P\{(c_j^{\mathsf T}Q\phi)^2\}=E_P\{\Lambda_j(X)\}>0\); more generally, \(c^{\mathsf T}(V_{\mathrm{raw}}-V_{\mathrm{eff}})c=E_P\{(c^{\mathsf T}Q\phi)^2\}\) for every \(c\in\mathbb R^K\).
\end{theorem}
\begin{proof}
Fix \(c_j\in\mathcal C_{\mathrm{adj}}\), abbreviate \(v=v_j\), \(b=b_v\), and put
\[
 e_0=E_P\{(N_Bv)^2\mid X\}. \tag{26}
\]
The augmented matrix is the conditional Gram matrix of \((h^{\mathsf T},N_Bv)^{\mathsf T}\), so adjoining its final coordinate increases rank by at most one.  If \(a\in\ker G\), then
\(E_P\{(a^{\mathsf T}h)^2\mid X\}=0\), hence \(a^{\mathsf T}h=0\) conditionally and \(a^{\mathsf T}b=0\).  Therefore \(b\in\operatorname{range}(G)\).  Multiplication on the left and right by the invertible block-triangular matrices that subtract \(G^+b\) from the final coordinate gives
\[
 \operatorname{rank}(G_j^{\mathrm{aug}})
 =\operatorname{rank}(G)
 \mathbf 1\{e_0-b^{\mathsf T}G^+b>0\}. \tag{27}
\]
The generalized Schur residual is
\[
\begin{aligned}
 e_0-b^{\mathsf T}G^+b
 &=E_P\!\left[\{N_Bv-h^{\mathsf T}G^+b\}^2\mid X\right]\\
 &=E_P\{(c_j^{\mathsf T}Q\phi)^2\mid X\}
 =\Lambda_j(X)\geq0. \tag{28}
\end{aligned}
\]
The first equality uses \(GG^+b=b\); the second uses linearity of \(Q\) and the definition of the cell vector \(v_j\).  Equations (27)--(28) prove equivalence of the positive-probability rank increase, \(P_X\{\Lambda_j>0\}>0\), and \(c_j^{\mathsf T}Q\phi\ne0\) in \(L^2(P)\).

Since \(Q\) is an orthogonal projection, \(Q\phi\) and \((I-Q)\phi=\phi^{\mathrm{eff}}\) are orthogonal coordinatewise.  Therefore, for every \(c\in\mathbb R^K\),
\[
\begin{aligned}
 c^{\mathsf T}(V_{\mathrm{raw}}-V_{\mathrm{eff}})c
 &=E_P\{(c^{\mathsf T}\phi)^2-(c^{\mathsf T}(I-Q)\phi)^2\}\\
 &=E_P\{(c^{\mathsf T}Q\phi)^2\}. \tag{29}
\end{aligned}
\]
For \(c=c_j\), iterating (28) over \(X\) gives the equality with \(E_P\{\Lambda_j(X)\}\), and it is strictly positive exactly under the observable rank-increase certificate.
\end{proof}
\par\smallskip\noindent \textit{Justification.} The conditional residual-energy diagnostic can be estimated from recovered cells and the same Gram solve, so strict gain for each adjacent exposure step is empirically certifiable rather than assumed. \textit{Gap.} Categorical product-bridge validity and the abstract projection formula do not provide an observable criterion for deciding which scientifically prespecified adjacent smoking-intensity contrasts receive a nonzero correction. \textit{Consumer.} It lets the maternal-smoking analyst report, for every neighboring pair of exposure bins, whether canonicalization yields a genuine efficiency gain.

\begin{lemma}[lem:6-score-stability]\label{lem:6-score-stability}
For every regular coherent foldwise fit, let \(\widehat M_A\) and \(\widehat M_Z\) denote the raw measurement fits and let \(\bar M_A\) and \(\bar M_Z\) denote their threshold-and-renormalize cleanings.  Every fitted bridge and every fitted object \(\widehat B,\widehat d,\widehat D,\widehat R,\widehat N,\widehat h,\widehat G,\widehat b_{(\cdot)},\widehat Q\) is formed from \((\bar M_A,\bar M_Z)\), and \(\widehat m_j(x)=\int_{\mathcal Y}y\widehat f_j(y\mid x)d\nu(y)\).  With \(r_A^c=r_A+\tau_{S,n}\) and \(r_Z^c=r_Z+\tau_{S,n}\),
\[
 \{E_P\|\widehat N-N_B\|_{\mathrm{op}}^2\}^{1/2}
 \leq C_{\mathrm{reg}}(r_D+r_B),
\]
\(r_R+r_N+r_h+r_G+r_b=O_P(R_\Sigma)\), and
\[
 \|\widehat Q\widehat U-QU\|_{L^2(P)}
 \leq C_{\mathrm{reg}}R_\Sigma.
\]
The full estimated efficient score \(\widehat U-\widehat Q\widehat U\) has the same linear \(L^2(P)\) stability bound relative to \(U-QU\).  In particular, the bridge perturbations use
\[
 \bar M_A^{-1}M_A-I_K=\bar M_A^{-1}(M_A-\bar M_A),
 \qquad
 \bar M_Z^{-1}M_Z-I_K=\bar M_Z^{-1}(M_Z-\bar M_Z),
\]
so the raw essential-supremum errors \(r_{A,\infty}\) and \(r_{Z,\infty}\) are screening quantities only, while \(r_B,\ldots,r_b\) are errors of the cleaned-law plug-ins.
\end{lemma}
\begin{proof}
Write \(\|V\|_2=\{E_P\|V\|^2\}^{1/2}\), using the Euclidean, Frobenius, or operator norm according to the type of \(V\).  Constants below depend only on the quantities allowed in \(C_{\mathrm{reg}}\).  Lemma \(\mathrm{lem{:}support\mbox{-}leakage\mbox{-}conditional\mbox{-}zero}\) gives
\[
 \|\bar M_A-M_A\|_2\leq C_{\mathrm{reg}}r_A^c,
 \qquad
 \|\bar M_Z-M_Z\|_2\leq C_{\mathrm{reg}}r_Z^c. \tag{30}
\]
These bounds do not require the screening event; that event is needed only for exact recovery of zero cells.

Put
\[
 H=B^{\mathsf T}D^{-1}B,\qquad
 \widehat H=\widehat B^{\mathsf T}\widehat D^{-1}\widehat B,
 \qquad \widehat B=\bar M_A\odot\bar M_Z. \tag{31}
\]
Observed-cell positivity gives \(\|D^{-1}\|_{\mathrm{op}}\leq\varepsilon^{-1}\), and the fitted half-margin gives \(\|\widehat D^{-1}\|_{\mathrm{op}}\leq2\varepsilon^{-1}\).  For \(v\in\mathbb R^K\), matricization yields
\[
 \operatorname{mat}(Bv)=M_A\operatorname{diag}(v)M_Z^{\mathsf T},
 \qquad
 \operatorname{mat}(\widehat Bv)=\bar M_A\operatorname{diag}(v)\bar M_Z^{\mathsf T}. \tag{32}
\]
The true and fitted singular-value margins imply
\[
 \|Bv\|_2\geq\varepsilon^2\|v\|_2,
 \qquad
 \|\widehat Bv\|_2\geq\frac{\varepsilon^2}{4}\|v\|_2. \tag{33}
\]
Because the diagonal entries of \(D\) and \(\widehat D\) are probabilities and hence at most one,
\[
 v^{\mathsf T}Hv\geq\varepsilon^4\|v\|_2^2,\qquad
 v^{\mathsf T}\widehat Hv\geq\frac{\varepsilon^4}{16}\|v\|_2^2.
\]
Consequently
\[
 \|H^{-1}\|_{\mathrm{op}}\leq\varepsilon^{-4},
 \qquad
 \|\widehat H^{-1}\|_{\mathrm{op}}\leq16\varepsilon^{-4}. \tag{34}
\]

The exact inverse identities
\[
 \widehat D^{-1}-D^{-1}=\widehat D^{-1}(D-\widehat D)D^{-1},
 \qquad
 \widehat H^{-1}-H^{-1}=\widehat H^{-1}(H-\widehat H)H^{-1} \tag{35}
\]
and product telescoping in \(B^{\mathsf T}D^{-1}B\) give
\[
 \|\widehat H-H\|_{\mathrm{op}}
 \leq C_{\mathrm{reg}}
 \{\|\widehat B-B\|_{\mathrm{op}}+\|\widehat D-D\|_{\mathrm{op}}\}. \tag{36}
\]
Applying (35)--(36) to the formulas for \(N_B\) and \(\widehat N\) yields pointwise
\[
 \|\widehat N-N_B\|_{\mathrm{op}}
 \leq C_{\mathrm{reg}}
 \{\|\widehat D-D\|_{\mathrm{op}}+\|\widehat B-B\|_{\mathrm{op}}\}. \tag{37}
\]
Taking the declared conditional \(L^2(P)\) norm proves the first displayed assertion of the lemma.

The cleaned inverse-feature bridges satisfy the exact identities
\[
 \bar M_A^{-1}M_A-I_K=\bar M_A^{-1}(M_A-\bar M_A),
 \qquad
 \bar M_Z^{-1}M_Z-I_K=\bar M_Z^{-1}(M_Z-\bar M_Z). \tag{38}
\]
The fitted singular-value half-margins and (30) therefore bound their conditional bridge-mean errors by \(C_{\mathrm{reg}}r_A^c\) and \(C_{\mathrm{reg}}r_Z^c\).  Since
\(\widehat m_j=\int_{\mathcal Y}y\widehat f_jd\nu\) and the fitted density is supported on \(\mathcal Y\), bounded outcome support gives \(|\widehat m_j|\leq L\).  Product telescoping in the bridge formula, latent positivity, and the fitted propensity half-margin give
\[
 \|\widehat U-U\|_2
 \leq C_{\mathrm{reg}}(r_m+r_\pi+r_A^c+r_Z^c)
 \leq C_{\mathrm{reg}}R_\Sigma. \tag{39}
\]

Lemma \(\mathrm{lem{:}data\mbox{-}driven\mbox{-}gram\mbox{-}threshold}\), with \(\tau_G=\varepsilon/2\), gives
\[
 \|\widehat G_{\tau_G}^{\dagger}-G^+\|_{L^2(P_X;\mathrm{op})}
 \leq26\varepsilon^{-2}r_G,\qquad
 \|\widehat G_{\tau_G}^{\dagger}\|_{\mathrm{op}}\leq2\varepsilon^{-1}. \tag{40}
\]
Every preceding map is evaluated at the cleaned product law.  Conditional expectation is an \(L^2\) contraction, and the inverse identities and finite product telescoping propagate the declared cleaned-law errors through \(\widehat B,\widehat d,\widehat D,\widehat R,\widehat N,\widehat c,\widehat h,\widehat G\), and \(\widehat b_{(\cdot)}\).  Thus
\[
 r_R+r_N+r_h+r_G+r_b=O_P(R_\Sigma). \tag{41}
\]
Neither \(r_{A,\infty}\) nor \(r_{Z,\infty}\) appears in (39)--(41); they enter only the screening ratio.

Let \(\widehat b_{\widehat U}\) be the right-hand side formed from the same cleaned fitted law and \(\widehat U\).  Add and subtract one factor at a time:
\[
\begin{aligned}
 &\widehat h^{\mathsf T}\widehat G_{\tau_G}^{\dagger}\widehat b_{\widehat U}
   -h^{\mathsf T}G^+b_U\\
 &=(\widehat h-h)^{\mathsf T}\widehat G_{\tau_G}^{\dagger}\widehat b_{\widehat U}
   +h^{\mathsf T}(\widehat G_{\tau_G}^{\dagger}-G^+)\widehat b_{\widehat U}
   +h^{\mathsf T}G^+(\widehat b_{\widehat U}-b_U).
\end{aligned} \tag{42}
\]
Conditional Cauchy--Schwarz, (40), and the margin bounds imply
\[
 \|\widehat h^{\mathsf T}\widehat G_{\tau_G}^{\dagger}\widehat b_{\widehat U}
   -h^{\mathsf T}G^+b_U\|_2
 \leq C_{\mathrm{reg}}(r_h+r_G+r_b)
 \leq C_{\mathrm{reg}}R_\Sigma. \tag{43}
\]
Likewise, (37) and (39) imply
\[
 \|\widehat N\widehat U-N_BU\|_2
 \leq C_{\mathrm{reg}}R_\Sigma. \tag{44}
\]
Subtracting (43) from (44) according to the definitions of \(\widehat Q\) and \(Q\) proves
\[
 \|\widehat Q\widehat U-QU\|_2\leq C_{\mathrm{reg}}R_\Sigma. \tag{45}
\]
Finally,
\[
(\widehat U-\widehat Q\widehat U)-(U-QU)
=(\widehat U-U)-(\widehat Q\widehat U-QU),
\]
so (39), (45), and the triangle inequality prove the same linear \(L^2(P)\) stability bound for the complete efficient score.
\end{proof}
\par\smallskip\noindent \textit{Justification.} The revised stability proof uses barred measurement inverses and propagates only cleaned-law errors through the posterior and Gram maps. \textit{Gap.} Existing bridge papers do not analyze perturbation of this support-cleaned nested projector. \textit{Consumer.} It turns a coherent cleaned factor fit into stable efficient scores and covariance estimates.

\begin{lemma}[lem:7-quadratic-bias]\label{lem:7-quadratic-bias}
At every two-sided pathwise-interior \(P\in\mathcal M_{K,\varepsilon}\), for every coherent foldwise fit whose raw measurement matrices are thresholded and renormalized before any bridge or projection object is formed, if \(\kappa_{\mathrm{supp},n}=o_P(1)\), then on an event whose probability tends to one,
\[
 \|E_P(\widehat U-\widehat Q\widehat U)-\psi\|
 \leq C_{\mathrm{reg}}\Delta_n,
 \qquad
 \Delta_n=r_m(r_\pi+r_A^c+r_Z^c)+r_A^cr_Z^c
 +(r_A^c+r_Z^c)R_\Sigma.
\]
Here \(\widehat m_j=\int_{\mathcal Y}y\widehat f_jd\nu\), and \(r_B,\ldots,r_b\) refer to the cleaned-law plug-ins.  In particular, the pathwise derivative of the mean bias is zero along a support-preserving two-sided perturbation of \(M_A\) alone and along a coherent perturbation of \(f\) alone.
\end{lemma}
\begin{proof}
Work coordinatewise and set \(k=QU\) and \(\widehat k=\widehat Q\widehat U\).  Lemma \(\mathrm{lem{:}support\mbox{-}leakage\mbox{-}conditional\mbox{-}zero}\) implies that \(\kappa_{\mathrm{supp},n}=o_P(1)\) produces an event \(\mathcal E_{\mathrm{clean},n}\) of probability tending to one on which
\[
 \|E_P\widehat k\|
 \leq C_{\mathrm{reg}}
 \left[(r_A^c+r_Z^c)\|\widehat k-k\|_2+r_A^cr_Z^c\right]. \tag{46}
\]
Lemma \(\mathrm{lem{:}6\mbox{-}score\mbox{-}stability}\) gives
\(\|\widehat k-k\|_2\leq C_{\mathrm{reg}}R_\Sigma\).  Hence
\[
 \|E_P\widehat k\|
 \leq C_{\mathrm{reg}}\{(r_A^c+r_Z^c)R_\Sigma+r_A^cr_Z^c\}. \tag{47}
\]

For \(j,\ell\in\mathcal J\), define the cleaned bridge-mean errors
\[
\begin{aligned}
 \delta^A_{j\ell}(X)
 &=E_P\{\bar a_j(A,X)\mid J=\ell,X\}-\mathbf 1\{j=\ell\},\\
 \delta^Z_{j\ell}(X)
 &=E_P\{\bar z_j(Z,X)\mid J=\ell,X\}-\mathbf 1\{j=\ell\}.
\end{aligned} \tag{48}
\]
The bridge definitions and true measurement laws give
\[
 \delta^A=\bar M_A^{-1}M_A-I_K=\bar M_A^{-1}(M_A-\bar M_A),
 \quad
 \delta^Z=\bar M_Z^{-1}M_Z-I_K=\bar M_Z^{-1}(M_Z-\bar M_Z). \tag{49}
\]
The fitted half-margins and cleaning bounds imply
\[
 \|\delta^A\|_{L^2(P_X)}\leq C_{\mathrm{reg}}r_A^c,\qquad
 \|\delta^Z\|_{L^2(P_X)}\leq C_{\mathrm{reg}}r_Z^c, \tag{50}
\]
and both matrices are uniformly bounded.

By three-view conditional independence, \(Y,A,Z\) separate conditionally on \((J=\ell,X)\), and
\(E_P(Y-\widehat m_j\mid J=\ell,X)=m_\ell-\widehat m_j\).  Expanding \(\widehat U_j\), conditioning on \((J,X)\), and summing over \(J\) gives the exact identity
\[
\begin{aligned}
 E_P(\widehat U_j\mid X)-m_j
={}&(\widehat m_j-m_j)
 \left[1-\frac{\pi_j}{\widehat\pi_j}
 (1+\delta^A_{jj})(1+\delta^Z_{jj})\right]\\
&+\widehat\pi_j^{-1}\sum_{\ell\ne j}\pi_\ell
 \delta^A_{j\ell}\delta^Z_{j\ell}(m_\ell-\widehat m_j). \tag{51}
\end{aligned}
\]
For \(\ell=j\), both bridge means are one plus their errors; for \(\ell\ne j\), both indicator terms vanish.  This proves (51) without approximation.

Latent positivity and the fitted propensity half-margin give
\(\pi_j^{-1}\leq\varepsilon^{-1}\) and \(\widehat\pi_j^{-1}\leq2\varepsilon^{-1}\).  Bounded support gives \(|m_\ell|\vee|\widehat m_j|\leq L\), and
\[
 \left|1-\frac{\pi_j}{\widehat\pi_j}\right|
 \leq2\varepsilon^{-1}|\widehat\pi_j-\pi_j|. \tag{52}
\]
Expanding the first bracket in (51), using uniform boundedness of the two \(\delta\)'s, and applying Cauchy--Schwarz with (50) bounds its expectation by
\(C_{\mathrm{reg}}r_m(r_\pi+r_A^c+r_Z^c)\).  Cauchy--Schwarz applied to the second line bounds its expectation by \(C_{\mathrm{reg}}r_A^cr_Z^c\).  Since \(K\) is fixed and \(E_Pm_j(X)=\psi_j\),
\[
 \|E_P\widehat U-\psi\|
 \leq C_{\mathrm{reg}}
 \{r_m(r_\pi+r_A^c+r_Z^c)+r_A^cr_Z^c\}. \tag{53}
\]
Using
\[
E_P(\widehat U-\widehat Q\widehat U)-\psi
=(E_P\widehat U-\psi)-E_P\widehat k
\]
and combining (47), (53), and the definition of \(\Delta_n\) proves the displayed bias bound on \(\mathcal E_{\mathrm{clean},n}\).

For the derivative assertions, let \(t\mapsto\eta_t\) be a support-preserving two-sided regular perturbation with \(\eta_0=\eta\), and evaluate score means under \(P=P_{\eta_0}\).  If only \(M_A\) varies, then \(m_t=m\), \(\pi_t=\pi\), and \(\delta_t^Z=0\).  Equation (51) gives
\[
 E_PU(\eta_t)-\psi=0. \tag{54}
\]
Support preservation makes the true-zero discrepancy identically zero.  The conditional-zero argument of the support-cleaning lemma, with the \(Z\)-error zero, gives
\[
 \|E_PQ(\eta_t)U(\eta_t)\|
 \leq C_{\mathrm{reg}}r_A(t)
 \|Q(\eta_t)U(\eta_t)-QU\|_2=O(t^2), \tag{55}
\]
because DQM differentiability gives \(r_A(t)=O(|t|)\), while the separated inverse and Gram-gap perturbation identities give score displacement \(O(|t|)\).  Equations (54)--(55) make the derivative at zero vanish.

If only \(f\) varies coherently, then \(M_A,M_Z,\pi\), and \(B\) are fixed, so
\(\delta_t^A=\delta_t^Z=0\).  Equation (51) again gives \(E_PU(\eta_t)=\psi\).  Projector algebra gives \(B^{\mathsf T}N_t=0\), and every column of \(h_t=N_tc_t\) belongs to \(\ker(B^{\mathsf T})\).  Therefore
\[
 B^{\mathsf T}\{Q_tU(\eta_t)\}=0. \tag{56}
\]
By three-view independence, the \(j\)-th coordinate of (56) is
\[
 E_P\{Q_tU(\eta_t)\mid Y,J=j,X\}
 =\sum_{a,z}B_{(a,z),j}\{Q_tU(\eta_t)\}_{az}(Y,X)=0.
\]
Iterated expectation gives \(E_PQ_tU(\eta_t)=0\) exactly.  Thus the mean bias is identically zero along the coherent \(f\)-only perturbation, proving the second derivative assertion.
\end{proof}
\par\smallskip\noindent \textit{Justification.} The corrected expansion uses \(\bar M_A^{-1}M_A-I\) and \(\bar M_Z^{-1}M_Z-I\), so coherent support cleaning leaves only products of cleaned errors. \textit{Gap.} The neighboring binary and categorical results do not cover this cleaned multilevel posterior-and-Gram correction. \textit{Consumer.} It provides the explicit second-order remainder used by the cross-fitted estimator.

\begin{theorem}[thm:8-efficient-estimation]\label{thm:8-efficient-estimation}
At every two-sided pathwise-interior \(P\in\mathcal M_{K,\varepsilon}\), let every training fold use the support-cleaned coherent fit and the data-driven Gram inverse at \(\tau_G=\varepsilon/2\).  Suppose, for every fold,
\[
 \kappa_{\mathrm{supp},n}=o_P(1),
 \qquad \tau_{S,n}=o_P(n^{-1/4}),
 \qquad R_\Sigma=o_P(1),
 \qquad \Delta_n=o_P(n^{-1/2}).
\]
Then
\[
 \sqrt n(\widehat\psi-\psi)
 =n^{-1/2}\sum_{i=1}^n\phi^{\mathrm{eff}}(O_i)+o_P(1),
 \qquad
 \sqrt n(\widehat\psi-\psi)\rightsquigarrow N(0,V_{\mathrm{eff}}),
\]
and \(\widehat V_{\mathrm{eff}}\to_PV_{\mathrm{eff}}\).  The raw matrices \(\widehat M_A,\widehat M_Z\) are used only to construct \(\bar M_A,\bar M_Z\) and the screening ratio; all bridges, posterior objects, projectors, and scores use the cleaned matrices.
\end{theorem}
\begin{proof}
For a training fold \(r\), define
\[
 S(O)=U(O)-QU(O)=\psi+\phi^{\mathrm{eff}}(O),\qquad
 \widehat S_r=\widehat U_r-\widehat Q_r\widehat U_r, \tag{57}
\]
where every hatted object is computed from the cleaned product law.  Let \(n_r\) be the held-out fold size.  Adding and subtracting foldwise conditional means gives
\[
\begin{aligned}
 \sqrt n(\widehat\psi-\psi)
={}&n^{-1/2}\sum_{i=1}^n\phi^{\mathrm{eff}}(O_i)\\
&+n^{-1/2}\sum_{r=1}^{F_{\mathrm{fold}}}
 \sum_{i:\,\mathrm{fold}(i)=r}
 [(\widehat S_r-S)(O_i)-P(\widehat S_r-S)]\\
&+n^{-1/2}\sum_{r=1}^{F_{\mathrm{fold}}}n_rP(\widehat S_r-\psi). \tag{58}
\end{aligned}
\]
Conditional on the training observations for fold \(r\), its held-out observations are independent of \(\widehat S_r\), and the conditional second moment of its centered contribution is at most
\[
 \frac{n_r}{n}\|\widehat S_r-S\|_2^2. \tag{59}
\]
Lemma \(\mathrm{lem{:}6\mbox{-}score\mbox{-}stability}\) and \(R_\Sigma=o_P(1)\) make (59) converge to zero in probability.  Conditional Markov inequality and fixed \(F_{\mathrm{fold}}\) make the middle line of (58) \(o_P(1)\).

The support-cleaning event has probability tending to one because
\(\kappa_{\mathrm{supp},n}=o_P(1)\).  On that event Lemma \(\mathrm{lem{:}7\mbox{-}quadratic\mbox{-}bias}\) gives, foldwise,
\[
 \left\|n^{-1/2}\sum_rn_rP(\widehat S_r-\psi)\right\|
 \leq\sqrt n\,C_{\mathrm{reg}}\max_r\Delta_{n,r}=o_P(1). \tag{60}
\]
The condition \(\tau_{S,n}=o_P(n^{-1/4})\) allows a training-fold-random threshold and places each cleaning increment
\(r_A^c-r_A=r_Z^c-r_Z=\tau_{S,n}\) in the stated product-rate window; the needed bias conclusion is explicitly imposed through
\(\Delta_n=o_P(n^{-1/2})\).  Equations (58)--(60) prove the asymptotic linear expansion.

The true regularity margins, bounded outcome, fixed \(K\), and the finite-projector and Gram-gap bounds make \(\phi^{\mathrm{eff}}\) bounded and centered.  For fixed \(t\in\mathbb R^K\), Taylor expansion of its characteristic function, with a bounded third-order remainder, gives
\[
 E_P\exp\{it^{\mathsf T}\phi^{\mathrm{eff}}/\sqrt n\}
 =1-\frac{t^{\mathsf T}V_{\mathrm{eff}}t}{2n}+o(n^{-1}). \tag{61}
\]
By Assumption \(\mathrm{ass{:}iid}\), raising (61) to the \(n\)-th power yields
\(\exp\{-t^{\mathsf T}V_{\mathrm{eff}}t/2\}\).  The Cramér--Wold device therefore gives
\[
n^{-1/2}\sum_{i=1}^n\phi^{\mathrm{eff}}(O_i)
\rightsquigarrow N(0,V_{\mathrm{eff}}),
\]
and the expansion proves the asserted limit for \(\widehat\psi\).

For covariance consistency, score stability gives \(\widehat S_r\to S\) in \(L^2(P)\) for every fold.  The fitted half-margins, spectral cutoff, and bounded outcome make these scores uniformly bounded.  Conditional on the training data, a direct second-moment calculation gives
\[
 \left\|P_{n,r}(\widehat S_r\widehat S_r^{\mathsf T})
 -P(\widehat S_r\widehat S_r^{\mathsf T})\right\|_{\mathrm F}=o_P(1). \tag{62}
\]
Cauchy--Schwarz gives
\[
 \|P(\widehat S_r\widehat S_r^{\mathsf T})-P(SS^{\mathsf T})\|_{\mathrm F}
 \leq\|\widehat S_r-S\|_2(\|\widehat S_r\|_2+\|S\|_2)=o_P(1). \tag{63}
\]
Summing the fixed number of folds and using \(\widehat\psi\to_P\psi\) yields
\[
 \widehat V_{\mathrm{eff}}
 \to_P E_P\{(S-\psi)(S-\psi)^{\mathsf T}\}=V_{\mathrm{eff}}.
\]
\end{proof}
\par\smallskip\noindent \textit{Justification.} The attaining theorem now exposes the exact screening, random-threshold, stability, and quadratic-bias conditions required by the non-oracle cleaned estimator. \textit{Gap.} Prior hidden-treatment limit theorems do not combine the exact multilevel projector with support cleaning and a data-driven Gram rank. \textit{Consumer.} It yields efficient standard errors for latent maternal-smoking exposure means when the declared rate conditions are verified.

\begin{theorem}[thm:9-joint-inference]\label{thm:9-joint-inference}
For fixed \(q\), when \(\widehat\psi\) and \(\widehat\Sigma_C\) use the corrected coherent projector in Theorem 8, \(\sqrt n(C\widehat\psi-C\psi)\rightsquigarrow N(0,\Sigma_C)\), \(\widehat\Sigma_C\to_P\Sigma_C\), and the Gaussian maximum-\(t\) critical value computed from the correlation matrix of \(\widehat\Sigma_C\) gives simultaneous coverage converging to \(1-\alpha\) for all \(q\) coordinates of \(\theta=C\psi\).
\end{theorem}
\begin{proof}
Multiplying the expansion in Theorem \(\mathrm{thm{:}8\mbox{-}efficient\mbox{-}estimation}\) by the fixed matrix \(C\) gives
\[
\begin{aligned}
 \sqrt n(C\widehat\psi-C\psi)
 &=n^{-1/2}\sum_{i=1}^nC\phi^{\mathrm{eff}}(O_i)+o_P(1)\\
 &\rightsquigarrow N(0,CV_{\mathrm{eff}}C^{\mathsf T})
 =N(0,\Sigma_C). \tag{64}
\end{aligned}
\]
Consistency of \(\widehat V_{\mathrm{eff}}\) and continuity of finite matrix multiplication give
\(\widehat\Sigma_C\to_P\Sigma_C\).  Assumption \(\mathrm{ass{:}contrast\mbox{-}information}\) makes every diagonal entry of \(\Sigma_C\) positive, so standardization is continuous and the estimated correlation matrix \(\widehat R_C\) converges in probability to the correlation matrix \(R_C\).

Because \(q\) is fixed and \(R_C\) is positive definite, a Gaussian vector with covariance \(R_C\) has a continuous positive density.  The boundary of every rectangle \([-t,t]^q\) has probability zero, and every nonempty shell between two such rectangles has positive probability.  Thus the distribution function of \(\max_{r\leq q}|Z_r|\), \(Z\sim N(0,R_C)\), is continuous and strictly increasing between zero and one.  If deterministic correlation matrices \(R_m\to R_C\), their Gaussian densities converge pointwise; domination on compact sets and a uniform Gaussian tail bound imply convergence of rectangle probabilities.  The corresponding \((1-\alpha)\)-quantiles therefore converge.  Applying this deterministic conclusion along every almost-surely convergent subsequence of \(\widehat R_C\) proves convergence in probability of the estimated Gaussian maximum-\(t\) critical value.

Combining critical-value convergence with (64) and consistent coordinate standard errors yields
\[
 P\!\left(
 |\sqrt n(C\widehat\psi-C\psi)_r|
 \leq\widehat c_{1-\alpha}\sqrt{\widehat\Sigma_{C,rr}}
 \text{ for all }r\leq q\right)\longrightarrow1-\alpha. \tag{65}
\]
Assumption \(\mathrm{ass{:}contrast\mbox{-}centering}\) makes every row of \(C\) a treatment contrast; full row rank and positive information exclude redundant or zero-variance reported coordinates.
\end{proof}
\par\smallskip\noindent \textit{Justification.} The application compares several adjacent smoking levels, so scalar pointwise intervals do not answer the scientific question. \textit{Gap.} The closest hidden-treatment theory is centered on a binary contrast and does not provide fixed-dimensional joint maximum-\(t\) inference for ordered multilevel effects. \textit{Consumer.} It supports one familywise-valid assessment of all declared six-bin maternal-smoking contrasts.

\begin{proposition}[prop:10-binary-reduction]\label{prop:10-binary-reduction}
In the binary restriction obtained by replacing \(K\geq3\) with \(J^{\mathrm{bin}}\in\{0,1\}\), use \(R=D^{-1}BW\) and \(N_B=I-D^{-1}B(B^{\mathsf T}D^{-1}B)^{+}B^{\mathsf T}\). The resulting cell-and-Gram operator \(Q\) reduces to the Zhou--Tchetgen Tchetgen continuous-outcome orthocomplement generated by \(B_0h_1(Y,X)+B_1h_2(Y,X)\) with \(E_P(h_r\mid J^{\mathrm{bin}},X)=0\). When \(Y\) is also binary, \(Q\phi=0\) and the raw product-bridge influence function is canonical. This restriction is a comparator sanity check and is not part of the proposal's \(K\geq3\) theorem scope.
\end{proposition}
\begin{proof}
In the binary restriction write \(\mu_{Aj}=E(A\mid J^{\mathrm{bin}}=j,X)\) and \(\mu_{Zj}=E(Z\mid J^{\mathrm{bin}}=j,X)\).  At fixed \((y,x)\), the two cell vectors
\[
B_0=(A-\mu_{A0})(Z-\mu_{Z1}),\qquad
B_1=(A-\mu_{A1})(Z-\mu_{Z0})
\]
satisfy \(B^{\mathsf T}B_0=B^{\mathsf T}B_1=0\): for each latent class, one of the two centered factors has conditional mean zero, and conditional independence factors the expectation.  The binary \(B\) has rank two, so \(\ker(B^{\mathsf T})\) has dimension two.  Measurement separation makes \(B_0,B_1\) linearly independent, hence every outcome-null cell function is uniquely
\[
k=B_0h_0(Y,X)+B_1h_1(Y,X). \tag{45}
\]
The restrictions conditional on \((A,J,X)\) applied successively at \(J=0,1\) imply \(E(h_0\mid J=0,X)=E(h_1\mid J=1,X)=0\); the restrictions conditional on \((Z,J,X)\) imply the two crossed equations.  Thus
\[
E(h_r\mid J^{\mathrm{bin}},X)=0,\qquad r=0,1. \tag{46}
\]
Conversely (46) makes all three conditional restrictions vanish.  The first \(D\)-orthogonal projection \(N_B\) removes the two-dimensional column space of \(B\); residualizing the four binary measurement representers and solving their conditional Gram system removes exactly the remaining violations of (46).  The same normal-equation calculation as (23)--(24), now wholly inside the binary restriction, therefore shows that \(Q\) is the orthogonal projection onto the space generated by (45)--(46).  This proves the Zhou--Tchetgen Tchetgen continuous binary reduction without invoking the \(K\geq3\) theorem.

If \(Y\) is binary, each \(h_r(Y,X)\) has two cell values.  The two equations in (46) have coefficient matrix formed by the two distinct class-conditional Bernoulli laws; the binary outcome-feature rank makes it nonsingular.  Hence \(h_0=h_1=0\), \(\operatorname{range}Q=\{0\}\), and \(Q\phi=0\).  This is only a restriction check and changes neither the \(K\ge3\) assumptions nor theorem scope.
\end{proof}
\par\smallskip\noindent \textit{Justification.} The reduction is a mandatory sanity check against the known binary locally nonparametric special case. \textit{Gap.} This is not a novelty claim; it verifies that the multilevel operator nests the published binary theory exactly. \textit{Consumer.} It permits direct validation of software and algebra against the published binary model before analyzing six exposure levels.

\begin{proposition}[prop:11-continuous-witness]\label{prop:11-continuous-witness}
The explicit non-binary law \(P^{\star}\) belongs to \(\mathcal M_{3,0.003}\). At each outcome bin its posterior weights are positive and nonconstant, \(R=D^{-1}BW\), and the exact first elimination is \(N_B=I_9-D^{-1}B(B^{\mathsf T}D^{-1}B)^{-1}B^{\mathsf T}\). The resulting operator is evaluable from nine \((A,Z)\) cells and one conditional Gram matrix of dimension at most eighteen; \(\operatorname{rank}(G)=12\), \(Q\phi\neq0\), \(\phi^{\mathrm{eff}}\in\mathcal T\), and \(V_{\mathrm{raw}}-V_{\mathrm{eff}}\succ0\).
\end{proposition}
\begin{proof}
All factors defining \(P^\star\) are rational and constant in \(X\).  On outcome bin \([b,b+1)\), the conditional mean is \(b+1/2\), so direct integration gives
\[
 m=\left(\frac9{10},\frac32,\frac{21}{10}\right)^{\mathsf T}. \tag{66}
\]
Every density entry lies in \([1/10,7/10]\), proving common support and the density envelope for \(\varepsilon=3/1000\); also \(|Y|\leq3\).  The three symmetric factor matrices have singular values
\[
 \sigma(M_A)=(1,7/10,7/10),\quad
 \sigma(M_Z)=(1,2/5,2/5),\quad
 \sigma(H)=(1,3/5,2/5). \tag{67}
\]
Here \(H=\Gamma_Y\), because \(h_0\) is the bin-indicator vector.  The modal gap of \(M_A\) is \(8/10-1/10=7/10\), and \(\min_j\pi_j=3/10\).

For bin \(b\), exact rational arithmetic gives
\[
 B_{(a,z),j}=M_{A,aj}M_{Z,zj},\quad
 w_{bj}=\frac{\pi_jH_{bj}}{\sum_\ell\pi_\ell H_{b\ell}},\quad
 D_b=\operatorname{diag}(Bw_b),\quad
 R_b=D_b^{-1}B\operatorname{diag}(w_b). \tag{68}
\]
The three posterior vectors before observing \((A,Z)\) are
\[
 w_0=(21/32,1/4,3/32)^{\mathsf T},\quad
 w_1=(1/6,2/3,1/6)^{\mathsf T},\quad
 w_2=(3/32,1/4,21/32)^{\mathsf T}, \tag{69}
\]
so each is positive and nonconstant.  Equation (67) and Lemma \(\mathrm{lem{:}3a\mbox{-}projector\mbox{-}algebra}\) make \(B\) full-column-rank and give
\[
 N_b=I_9-D_b^{-1}B(B^{\mathsf T}D_b^{-1}B)^{-1}B^{\mathsf T}. \tag{70}
\]
The smallest conditional cell probability is
\[
 \min_{b,a,z}(Bw_b)_{az}=\frac{69}{1600}>\frac3{1000}. \tag{71}
\]
Exact fraction-free \(LDL^{\mathsf T}\) elimination of
\(R_b^{\mathsf T}D_bR_b-(3/1000)I_3\) gives componentwise strict rational lower bounds for its pivot triples
\[
 (d_{01},d_{02},d_{03})\succ(547,149,35)/1000,\quad
 (d_{11},d_{12},d_{13})\succ(90,525,83)/1000, \tag{72}
\]
and
\[
 (d_{21},d_{22},d_{23})\succ(39,151,495)/1000. \tag{73}
\]
Each comparison follows by cross-multiplication of the rational entries in (68), so every pivot is positive.  Hence every posterior Gram eigenvalue exceeds \(3/1000\).

Construct \(N_b,h_b\) from (68) and integrate over the three unit-length bins.  For each class \(j\),
\(\sum_a h_{(j,a)}=N_bR_{b,\cdot j}=0\), and similarly
\(\sum_z h_{(j,z)}=0\).  Thus \(\operatorname{rank}(G)\leq2K(K-1)=12\).  Exact rational row reduction has pivot indices
\[
 I=\{1,2,4,5,7,8,10,11,13,14,16,17\}, \tag{74}
\]
and the corresponding principal minor is
\[
 \det(G_{I,I})=
 \frac{2^{29}3^{14}7^{10}\cdot47\cdot2887\cdot85346231\cdot258722383}
 {5^{40}\cdot2917\cdot536447^2\cdot138071\cdot21181156279^2}>0. \tag{75}
\]
Therefore \(\operatorname{rank}(G)=12\), and the six block-constant vectors span \(\ker(G)\).  Let \(S_0\in\mathbb R^{18\times12}\) be block diagonal with six copies of
\[
 S_{\mathrm{blk}}=
 \begin{pmatrix}1&0\\0&1\\-1&-1\end{pmatrix}. \tag{76}
\]
Its range is the Euclidean orthocomplement of that kernel.  Exact fraction-free elimination of
\[
 S_0^{\mathsf T}GS_0-\frac3{1000}S_0^{\mathsf T}S_0 \tag{77}
\]
gives twelve \(LDL^{\mathsf T}\) pivots with componentwise strict rational lower bounds
\[
 (23,4,8,23,20,3,81,27,37,29,72,20)/1000. \tag{78}
\]
Every comparison follows by cross-multiplication of the rational matrices in (68).  Hence (77) is positive definite.  If \(0\ne x\in\ker(G)^\perp\), write \(x=S_0a\); then
\[
 x^{\mathsf T}Gx-\frac3{1000}\|x\|_2^2
 =a^{\mathsf T}\left(S_0^{\mathsf T}GS_0-\frac3{1000}S_0^{\mathsf T}S_0\right)a>0. \tag{79}
\]
Thus every nonzero eigenvalue of \(G\) exceeds \(3/1000\), verifying both projection-Gram conditions.

The construction of \(Y^{\mathrm{pot}}\) makes \(Y(j)\perp J\), and \(Y=Y(J)\) gives consistency.  Mutual conditional independence of \(A,Z,Y^{\mathrm{pot}}\) given \(J\) gives the three-view condition.  Equations (66)--(79) verify every member property of \(\mathcal M_{3,0.003}\).  All inequalities are strict.  Bounded factorwise exponential tilts preserve them for sufficiently small parameters of either sign.  The structural upper bound \(\operatorname{rank}(G)\leq12\), together with the positive twelve-dimensional minor (75), keeps the rank equal to twelve locally.  Thus \(P^\star\) is a two-sided pathwise-interior point.

On each bin, \(R,D,N,h\) are constant and \(\phi\) is affine in \(u=y-b\in[0,1]\).  Hence every covariance is an exact rational expression using
\[
 \int_0^1(1,u,u^2)\,du=(1,1/2,1/3). \tag{80}
\]
The resulting raw covariance is
\[
 V_{\mathrm{raw}}=
 \begin{pmatrix}
 67699/8820&2927/7840&5689/13230\\
 2927/7840&1357/294&2927/7840\\
 5689/13230&2927/7840&67699/8820
 \end{pmatrix}. \tag{81}
\]
Writing \(K_0=Q\phi\), orthogonality of \(Q\) gives
\[
 V_{\mathrm{raw}}-V_{\mathrm{eff}}=E_P(K_0K_0^{\mathsf T}). \tag{82}
\]
Exact rational elimination places the three leading principal minors of the matrix in (82), in order, in
\[
 (1114/1000,1115/1000),\qquad
 (1219/1000,1220/1000),\qquad
 (1199/1000,1200/1000). \tag{83}
\]
These bounds are obtained by cross-multiplication of the rational entries generated by (68) and (80).  Gaussian elimination expresses its quadratic form as a sum of three squares with pivots equal to the first minor and the two successive ratios of leading minors.  Equation (83) makes all three pivots positive, so \(V_{\mathrm{raw}}-V_{\mathrm{eff}}\succ0\), and in particular \(Q\phi\ne0\).  Theorems \(\mathrm{thm{:}2\mbox{-}tangent\mbox{-}completeness}\) and \(\mathrm{thm{:}4\mbox{-}canonical\mbox{-}gradient}\) apply at this interior witness and give \(\phi^{\mathrm{eff}}\in\mathcal T\).  Finally, (68), (74), and the definition of \(Q\) show directly that computation uses nine cells and one Gram matrix of dimension at most eighteen.
\end{proof}
\par\smallskip\noindent \textit{Justification.} An exact continuous witness shows that strict gain is not a hand-tuned finite-table artifact and persists on a regular neighborhood. \textit{Gap.} The presolve calculation is numerical and quadrature-based; the proposition requires an analytic continuous score-membership and covariance calculation. \textit{Consumer.} It provides a reproducible unit test and a concrete design point for calibrating a multilevel hidden-exposure study.

\begin{openendedquestion}[oeq:2-coherent-fit-rates]\label{oeq:2-coherent-fit-rates}
On \(\mathcal M^{\beta}_{K,\varepsilon}(B_{\mathrm H})\) with \(\beta>(d_X+1)/2\), can the handle \(\mathfrak F\), using tensor-product B-splines of total dimension \(s_n\asymp n^{(d_X+1)/(2\beta+d_X+1)}\), be completed by explicit spline boundary conventions, a constrained product-likelihood parameterization, and a local likelihood-curvature and empirical-process argument that prove all of the following for one raw fit?  First,
\[
 r_\pi+r_A+r_Z+r_f=O_P(\rho_n),
 \qquad \rho_n=n^{-\beta/(2\beta+d_X+1)}.
\]
Second, with \(\zeta_n=\rho_n\sqrt{\log n}\),
\[
 \|\widehat\pi-\pi\|_\infty+
 \|\widehat M_A-M_A\|_\infty+
 \|\widehat M_Z-M_Z\|_\infty+
 \|\widehat f-f\|_\infty=O_P(\zeta_n).
\]
Third, the empirical minimum modal-gap rule over the \(K!\) common permutations selects the population labeling with probability tending to one.  Fourth, after threshold-and-renormalize cleaning at a training-fold threshold satisfying \(\zeta_n=o(\tau_{S,n})\) and \(\tau_{S,n}=o(n^{-1/4})\), the fitted propensity, cleaned measurement matrices, cell probabilities, and posterior Gram matrix satisfy the declared computable half-margin restrictions.  The cleaned support-rate-window lemma proves that such a threshold is compatible with the stated smoothness and yields exact true-zero recovery.  What remains open is the likelihood-curvature, spline approximation, localization, and uniform-rate proof for this specifically prescribed global product-likelihood fit; no conclusion about \(R_\Sigma\) or \(\Delta_n\) is asserted merely from a list of coordinatewise \(L^2\) rates.
\end{openendedquestion}
\par\smallskip\noindent \textit{Justification.} The residual question is now falsifiable at the missing likelihood-curvature, spline, primitive L2, uniform-rate, label, and half-margin steps. \textit{Gap.} The frozen smoothness assumptions and adjacent literature do not prove those facts for the specified global product likelihood. \textit{Consumer.} Resolving it would turn the conditional efficient-inference theorem into a complete implementation guarantee on the smooth subclass.

\begin{theorem}[thm:differential-score-lifting]\label{thm:differential-score-lifting}
For every \(P\in\mathcal M_{K,\varepsilon}\), the corrected handle \(\mathfrak I\) admits a measurable uniformly bounded right inverse on every observed-data DQM score: if \(s\in L^2_0(P)\) is the score of an observed path through \(P\), there are factorwise centered scores \(s_X,s_\pi,s_Y,s_A,s_Z\) such that
\[
s=E_P\{s_X(X)+s_\pi(J,X)+s_Y(Y,J,X)+s_A(A,J,X)+s_Z(Z,J,X)\mid O\},
\]
and the sum of their factor \(L^2\)-norms is at most \(C_{\mathrm{reg}}\|s\|_{L^2(P)}\).  The recovery is measurable in \(X\), uses the augmented finite-feature inverse to fix the differentiated factorization, then uses the conditional \(G(X)^+\) solve for the measurement scores and the pointwise \(\{R^{\mathsf T}DR\}^{-1}\) solve for the remaining mixture and class-density scores.  Class-density scores are support preserving.
\end{theorem}
\begin{proof}
Let \(s\in L^2_0(P)\) be the score of an observed-data DQM path.  Define
\[
 s_X(X)=E_P(s\mid X),\qquad s_c=s-s_X. \tag{84}
\]
Then \(E_P(s_c\mid X)=0\), while conditional Jensen gives
\(\|s_X\|_2\leq\|s\|_2\) and \(\|s_c\|_2\leq2\|s\|_2\).  For every bounded conditional moment \(H(O)\), DQM and Cauchy--Schwarz give its conditional derivative as \(E_P(Hs_c\mid X)\), bounded in absolute value by
\[
 \{E_P(H^2\mid X)\}^{1/2}\{E_P(s_c^2\mid X)\}^{1/2}. \tag{85}
\]
Apply Lemma \(\mathrm{lem{:}2a\mbox{-}constant\mbox{-}x\mbox{-}feature\mbox{-}inverse}\) pointwise to the derivatives of the normalized finite moments.  Its uniform right-inverse bound and (85) yield measurable derivatives of \(\pi,M_A,M_Z,\Gamma_Y\) whose conditional Euclidean norm is at most
\[
 C_{\mathrm{reg}}\{E_P(s_c^2\mid X)\}^{1/2}. \tag{86}
\]
The modal permutation is locally constant by its positive gap, so it has zero derivative.  Measurability follows because conditional moments have measurable versions on the standard Borel space \(\mathcal X\), and the finite inverse is Borel on the separated fixed-rank stratum.

Differentiate the identified density equation
\[
 p_A(y\mid x)=M_A(x)\Pi(x)f(y\mid x),\qquad
 \Pi=\operatorname{diag}(\pi), \tag{87}
\]
in conditional \(L^1(\nu)\).  Observed DQM implies the \(L^1\) derivative by Cauchy--Schwarz, and the uniformly bounded inverses in (87) give
\[
 \dot f=\Pi^{-1}M_A^{-1}
 \{\dot p_A-\dot M_A\Pi f-M_A\dot\Pi f\}. \tag{88}
\]
Its columns integrate to zero by normalization.  Common support permits the preliminary definition \(u_{Y,j}=\dot f_j/f_j\), \(f_jd\nu\)-almost everywhere.  Define \(u_{\pi,j}=\dot\pi_j/\pi_j\), and define each categorical factor score by derivative divided by a positive atom probability.  At a zero atom, a two-sided nonnegative probability path has derivative zero, and the score is set to zero.  Differentiating the product factorization now gives, almost everywhere,
\[
 s_c(y,a,z,x)=\sum_{j=1}^KR_{(a,z),j}(y,x)
 \{u_{\pi,j}(x)+u_{Y,j}(y,x)+u_{A,aj}(x)+u_{Z,zj}(x)\}. \tag{89}
\]
Equation (89) is used only to locate the finite cell component; the construction below replaces these preliminary scores by a bounded-norm representation.

Let \(u\in\mathbb R^{2K^2}\) collect the preliminary measurement-score coordinates.  Terms in (89) depending only on \((Y,J,X)\) have cell vectors in
\(\operatorname{range}(R)=\ker(N_B)\) by Lemma \(\mathrm{lem{:}3a\mbox{-}projector\mbox{-}algebra}\), whereas the measurement contribution is \(c^{\mathsf T}u\).  Consequently
\[
 N_Bs_c=h^{\mathsf T}u. \tag{90}
\]
Taking conditional inner products with \(h\) gives
\[
 b_s(X):=E_P(hN_Bs_c\mid X)=G(X)u. \tag{91}
\]
If \(v\in\ker G(X)\), then \(E_P\{(v^{\mathsf T}h)^2\mid X\}=0\), so \(v^{\mathsf T}h=0\) conditionally and \(v^{\mathsf T}b_s=0\).  Thus \(b_s\in\operatorname{range}(G)\).  Set \(u_0=G^+b_s\).  Since (91) has a solution, \(G(u_0-u)=0\), which implies \(h^{\mathsf T}(u_0-u)=0\); hence
\[
 h^{\mathsf T}u_0=N_Bs_c. \tag{92}
\]
The projection-Gram gap and least-squares projection identity give
\[
\begin{aligned}
 \|u_0(X)\|_2^2
 &=b_s^{\mathsf T}G^{+2}b_s
 \leq\varepsilon^{-1}b_s^{\mathsf T}G^+b_s\\
 &\leq\varepsilon^{-1}E_P\{(N_Bs_c)^2\mid X\}. \tag{93}
\end{aligned}
\]

For each class block, subtract from the \(A\)-coordinates of \(u_0\) their \(M_A(\cdot\mid j,X)\)-mean, and do the analogous operation for \(Z\).  A blockwise constant has cell image proportional to \(R_{\cdot j}\), which \(N_B\) annihilates.  Thus recentering preserves (92).  Since these are probability-weighted centering maps and \(K\) is fixed, the operation increases the Euclidean bound in (93) by at most a constant depending on \(K\).  Denote the centered coordinates by \((s_A,s_Z)\), and let \(c^{\mathsf T}u\) denote their observed cell contribution.

Set
\[
 s_0=s_c-c^{\mathsf T}u. \tag{94}
\]
Equation (92) implies \(N_Bs_0=0\).  Lemma 3A gives
\(\ker(N_B)=\operatorname{range}(R)\), so pointwise
\[
 t(y,x)=\{R^{\mathsf T}DR\}^{-1}R^{\mathsf T}D
 s_0(y,\cdot,x),\qquad s_0=Rt. \tag{95}
\]
The inverse exists by Assumption \(\mathrm{ass{:}posterior\mbox{-}gram\mbox{-}gap}\).  Since
\(s_0^{\mathsf T}Ds_0=t^{\mathsf T}R^{\mathsf T}DRt\) and \(0\leq w_j\leq1\),
\[
 \sum_{j=1}^Kw_j(y,x)t_j(y,x)^2
 \leq\|t(y,x)\|_2^2
 \leq\varepsilon^{-1}s_0(y,\cdot,x)^{\mathsf T}D(y,x)s_0(y,\cdot,x). \tag{96}
\]
Define
\[
 s_{\pi,j}(x)=\int_{\mathcal Y}f_j(y\mid x)t_j(y,x)d\nu(y),
 \qquad
 s_{Y,j}(y,x)=t_j(y,x)-s_{\pi,j}(x). \tag{97}
\]
Then \(E_P(s_Y\mid J,X)=0\).  The recentered measurement contribution has conditional mean zero given \(X\), and \(E_P(s_c\mid X)=0\).  Integrating \(s_0=Rt\) and using \(gw_j=\pi_jf_j\) gives
\[
 0=E_P(s_0\mid X)=\sum_{j=1}^K\pi_j(X)s_{\pi,j}(X), \tag{98}
\]
which is the propensity-score centering condition.  Equations (94)--(98) give the exact conditional-expectation representation in the theorem.

Integrate (96) against \(g(y\mid x)d\nu(y)\).  Since \(gw_j=\pi_jf_j\), the left side is the conditional latent \(L^2\) norm of \(t\).  Jensen applied to (97) bounds the joint norms of \(s_\pi\) and \(s_Y\) by a constant times this norm.  The pointwise \(D\)-orthogonal projector is a contraction in \(D\)-norm, while \(d_{az}\geq\varepsilon\) converts this into a uniform Euclidean bound for \(N_B\).  Combining these facts with (93)--(96), the finite-dimensional bound on the centered measurement contribution, and (84), gives
\[
 \|s_X\|_2+\|s_\pi\|_2+\|s_Y\|_2+\|s_A\|_2+\|s_Z\|_2
 \leq C_{\mathrm{reg}}\|s\|_2. \tag{99}
\]
All conditional expectations admit measurable versions, and every finite inverse is Borel on the declared separated fixed-rank strata, so the construction is measurable in \(X\).  Finally, common support makes (97) an \(f_jd\nu\)-almost-everywhere density score on the original support; it assigns no mass outside that support.  Bounded truncation followed by \(f_j\)-recentering converges in the class-factor \(L^2\) norm and yields support-preserving two-sided exponential tilts wherever paths are required in the tangent closure.  This proves the measurable uniformly bounded lift and support preservation.
\end{proof}

\begin{lemma}[lem:bian-categorical-product-bridge]\label{lem:bian-categorical-product-bridge}
\texttt{Appendix A.8 of Bian,\allowbreak{} Zhou,\allowbreak{} and Cui defines categorical inverse features,\allowbreak{} proves that their conditional expectations are latent-\allowbreak{}category indicators,\allowbreak{} forms pair-\allowbreak{}product bridges,\allowbreak{} and inserts them into a valid categorical-\allowbreak{}action influence function.}
\end{lemma}

\begin{lemma}[lem:ztt-partition-density]\label{lem:ztt-partition-density}
In the binary continuous-outcome construction of Supplement S8.2, the partition basis vector \(V^\ell\) is dense in the observed tangent-space orthocomplement as the equal-width partition resolution tends to infinity, on subsequences for which the displayed denominators and Gram inverses are well defined.
\end{lemma}

\begin{theorem}[thm:differential-inverse-resolution]\label{thm:differential-inverse-resolution}
The answer to the differential-inverse question is affirmative.  For every \(P\in\mathcal M_{K,\varepsilon}\), every observed-data DQM score admits a measurable conditional-\(X\), uniformly bounded lift to factorwise centered additive latent scores.  The lift uses \(\mathfrak I_0\) for the finite differentiated factors, the conditional \(G(X)^+\) solve for measurement scores, and the pointwise \(\{R^{\mathsf T}DR\}^{-1}\) solve for the residual propensity and support-preserving class-density scores.
\end{theorem}
\begin{proof}
Theorem \(\mathrm{thm{:}differential\mbox{-}score\mbox{-}lifting}\) applies to an arbitrary observed DQM score.  Its construction first recovers the \(X\)-marginal score and the finite differentiated factors measurably.  It then solves \(G(X)u_0=b_s(X)\) for centered measurement scores, with the projection-Gram gap giving the uniform \(L^2\) bound, and solves
\[
t=\{R^{\mathsf T}DR\}^{-1}R^{\mathsf T}Ds_0
\]
for the residual propensity and class-density components, with the posterior-Gram gap giving the remaining bound.  Equations (97)--(98) in that proof verify the outcome and propensity centering constraints, and its final truncation argument preserves the original class-density support.  These are precisely the measurable conditional-\(X\), uniformly bounded, support-preserving lift properties asserted here.  Therefore the differential-inverse question has the stated affirmative answer.
\end{proof}

\begin{lemma}[lem:data-driven-gram-threshold]\label{lem:data-driven-gram-threshold}
Let \(p=2K^2\), let \(G(x)\) be symmetric nonnegative definite, and suppose every nonzero eigenvalue of \(G(x)\) is at least \(\varepsilon\).  Form \(\widehat G_{\tau_G}^{\dagger}(x)\) from the fitted nonnegative-definite Gram matrix with \(\tau_G=\varepsilon/2\), without supplying a rank.  If \(\delta(x)=\|\widehat G(x)-G(x)\|_{\mathrm{op}}<\varepsilon/4\), then exactly \(\operatorname{rank}\{G(x)\}\) fitted eigenvalues exceed \(\tau_G\) and
\[
 \|\widehat G_{\tau_G}^{\dagger}(x)-G(x)^+\|_{\mathrm{op}}
 \leq 14\varepsilon^{-2}\delta(x).
\]
For arbitrary \(\delta(x)\), the integrated bound
\[
 \|\widehat G_{\tau_G}^{\dagger}-G^+\|_{L^2(P_X;\mathrm{op})}
 \leq 26\varepsilon^{-2}
 \|\widehat G-G\|_{L^2(P_X;\mathrm{op})}
\]
holds.  In particular, if \(r_{G,\infty}=\operatorname*{ess\,sup}_x\delta(x)=o_P(1)\), the data-driven selected rank equals the true rank for \(P_X\)-almost every \(x\) with probability tending to one.
\end{lemma}
\begin{proof}
Fix a realization of the training-fold fit and an \(x\) where the projection-Gram gap holds.  Suppress \(x\), put \(p=2K^2\), and order the eigenvalues as
\[
 \lambda_1\geq\cdots\geq\lambda_p\geq0,\qquad
 \widehat\lambda_1\geq\cdots\geq\widehat\lambda_p\geq0.
\]
Let \(\delta=\|\widehat G-G\|_{\mathrm{op}}\).  For every unit vector \(v\),
\[
 |v^{\mathsf T}(\widehat G-G)v|\leq\delta. \tag{100}
\]
For \(1\leq s\leq p\), intersect the bottom \((p-s+1)\)-dimensional eigenspace of \(G\) with the top \(s\)-dimensional eigenspace of \(\widehat G\).  The intersection contains a unit vector, and (100) gives
\(\widehat\lambda_s-\lambda_s\leq\delta\).  Interchanging the two matrices gives
\[
 |\widehat\lambda_s-\lambda_s|\leq\delta,\qquad 1\leq s\leq p. \tag{101}
\]

Write \(r=\operatorname{rank}(G)\).  The gap assumption gives
\[
 \lambda_s\geq\varepsilon\quad(1\leq s\leq r),\qquad
 \lambda_s=0\quad(r<s\leq p). \tag{102}
\]
If \(\delta<\varepsilon/4\), (101)--(102) and \(\tau_G=\varepsilon/2\) yield
\[
 \widehat\lambda_s>3\varepsilon/4>\tau_G\quad(s\leq r),\qquad
 \widehat\lambda_s<\varepsilon/4<\tau_G\quad(s>r). \tag{103}
\]
Thus exactly \(r\) eigenvalues are retained without supplying \(r\).

Define
\[
 \widetilde G=\sum_{s:\widehat\lambda_s>\tau_G}
 \widehat\lambda_s\widehat v_s\widehat v_s^{\mathsf T}.
\]
Then \(\widehat G_{\tau_G}^{\dagger}=\widetilde G^+\).  On the event (103), every discarded fitted eigenvalue is at most \(\delta\), so
\[
 \|\widetilde G-\widehat G\|_{\mathrm{op}}\leq\delta,\qquad
 \|\widetilde G-G\|_{\mathrm{op}}\leq2\delta, \tag{104}
\]
and
\[
 \|\widetilde G^+\|_{\mathrm{op}}\leq2\varepsilon^{-1},\qquad
 \|G^+\|_{\mathrm{op}}\leq\varepsilon^{-1}. \tag{105}
\]

For symmetric matrices \(H,A_0\), elementary Moore--Penrose algebra gives
\[
\begin{aligned}
 H^+-A_0^+
={}&-H^+(H-A_0)A_0^+
 +H^{+2}(H-A_0)(I-A_0A_0^+)\\
&+(I-H^+H)(H-A_0)A_0^{+2}. \tag{106}
\end{aligned}
\]
Indeed, using \(A_0(I-A_0A_0^+)=0\), \((I-H^+H)H=0\),
\(H^{+2}H=H^+\), and \(A_0A_0^{+2}=A_0^+\), the three terms reduce to
\(-P_HA_0^++H^+P_{A_0}\), \(H^+(I-P_{A_0})\), and
\(-(I-P_H)A_0^+\), whose sum is \(H^+-A_0^+\).

Apply (106) with \(H=\widetilde G\) and \(A_0=G\).  The complementary factors are orthogonal projections and have norm at most one.  Equations (104)--(105) bound the three terms by \(4\varepsilon^{-2}\delta\), \(8\varepsilon^{-2}\delta\), and \(2\varepsilon^{-2}\delta\).  Hence
\[
 \|\widehat G_{\tau_G}^{\dagger}-G^+\|_{\mathrm{op}}
 \leq14\varepsilon^{-2}\delta
 \quad\text{when }\delta<\varepsilon/4. \tag{107}
\]

Let \(\mathcal E=\{x:\delta(x)<\varepsilon/4\}\).  The threshold rule always gives
\(\|\widehat G_{\tau_G}^{\dagger}\|_{\mathrm{op}}\leq2\varepsilon^{-1}\), and the gap gives \(\|G^+\|_{\mathrm{op}}\leq\varepsilon^{-1}\).  Thus on \(\mathcal E^c\),
\[
 \|\widehat G_{\tau_G}^{\dagger}-G^+\|_{\mathrm{op}}
 \leq3\varepsilon^{-1}. \tag{108}
\]
Markov's inequality gives
\[
 P_X(\mathcal E^c)^{1/2}
 \leq4\varepsilon^{-1}\|\delta\|_{L^2(P_X)}. \tag{109}
\]
Using (107) on \(\mathcal E\), (108)--(109) on its complement, and the \(L^2(P_X)\) triangle inequality yields
\[
 \|\widehat G_{\tau_G}^{\dagger}-G^+\|_{L^2(P_X;\mathrm{op})}
 \leq26\varepsilon^{-2}
 \|\widehat G-G\|_{L^2(P_X;\mathrm{op})}. \tag{110}
\]
Finally, if \(r_{G,\infty}=o_P(1)\), then
\(\Pr\{r_{G,\infty}<\varepsilon/4\}\to1\).  On this event (103) holds for \(P_X\)-almost every \(x\), proving the rank-selection conclusion.
\end{proof}

\begin{lemma}[lem:support-leakage-conditional-zero]\label{lem:support-leakage-conditional-zero}
At every two-sided pathwise-interior \(P\in\mathcal M_{K,\varepsilon}\), consider a foldwise coherent fit.  Let \(\widehat M_A,\widehat M_Z\) be the raw measurement fits and \(\bar M_A,\bar M_Z\) their threshold-and-renormalize cleanings.  If \(\kappa_{\mathrm{supp},n}=o_P(1)\), then with probability tending to one every true-zero measurement cell is exactly zero in the corresponding cleaned matrix.  No lower bound on the positive entries of \(M_A\) or \(M_Z\) is required.  Moreover,
\[
 \|\bar M_A-M_A\|_\infty\leq C_{\mathrm{reg}}(r_{A,\infty}+\tau_{S,n}),
 \qquad
 \|\bar M_Z-M_Z\|_\infty\leq C_{\mathrm{reg}}(r_{Z,\infty}+\tau_{S,n}),
\]
and the analogous conditional \(L^2(P)\) bounds are \(C_{\mathrm{reg}}r_A^c\) and \(C_{\mathrm{reg}}r_Z^c\).  If \(k=QU\) and \(\widehat k=\widehat Q\widehat U\), coordinatewise, with all hatted derived objects formed from the cleaned matrices, then on the same event
\[
 \|E_P\widehat k\|
 \leq C_{\mathrm{reg}}\left[
 (r_A^c+r_Z^c)\|\widehat k-k\|_{L^2(P)}+r_A^cr_Z^c
 \right].
\]
Thus support cleaning removes the former first-order leakage term.
\end{lemma}
\begin{proof}
Fix one measurement view and suppress its subscript.  For a raw fitted probability column \(p=(p_1,\ldots,p_K)^{\mathsf T}\), define
\[
 q_a=p_a\mathbf 1\{p_a>\tau_{S,n}\},\qquad
 s=\sum_{b=1}^Kq_b,\qquad \bar p_a=q_a/s. \tag{111}
\]
Because \(\sum_ap_a=1\), the deleted mass \(t=1-s\) satisfies
\(0\leq t\leq K\tau_{S,n}\).  Since \(\tau_{S,n}\downarrow0\), eventually
\(\tau_{S,n}<1/(2K)\), so \(s\geq1/2\).  Hence
\[
 |\bar p_a-p_a|
 \leq|q_a-p_a|+q_a\frac{1-s}{s}
 \leq(2K+1)\tau_{S,n}. \tag{112}
\]
Applying (112) to every column and using the triangle inequality gives
\[
 \|\bar M_A-M_A\|_\infty
 \leq r_{A,\infty}+(2K+1)\tau_{S,n},\qquad
 \|\bar M_Z-M_Z\|_\infty
 \leq r_{Z,\infty}+(2K+1)\tau_{S,n}. \tag{113}
\]
The same pointwise inequality followed by the conditional \(L^2(P)\) norm gives
\[
 \|\bar M_A-M_A\|_2\leq C_{\mathrm{reg}}r_A^c,\qquad
 \|\bar M_Z-M_Z\|_2\leq C_{\mathrm{reg}}r_Z^c. \tag{114}
\]

Define
\[
 \mathcal E_{\mathrm{clean},n}
 =\{r_{A,\infty}+r_{Z,\infty}\leq\tau_{S,n}\}. \tag{115}
\]
If \(M_A(x)_{aj}=0\), then on (115),
\(\widehat M_A(x)_{aj}\leq\tau_{S,n}\), so (111) sets
\(\bar M_A(x)_{aj}=0\).  The same argument applies to \(Z\).  Since
\[
 \frac{r_{A,\infty}+r_{Z,\infty}}{\tau_{S,n}}
 =\kappa_{\mathrm{supp},n}=o_P(1), \tag{116}
\]
the probability of (115) tends to one.  No lower bound on a positive cell was used; a positive cell may also be deleted, and (112)--(114) account for that deletion.

Put \(\widehat B=\bar M_A\odot\bar M_Z\).  Fitted projector algebra gives
\(\widehat B^{\mathsf T}\widehat N=0\).  Both terms defining
\(\widehat k=\widehat Q\widehat U\) belong to \(\operatorname{range}(\widehat N)\), including under spectral thresholding because each coordinate of \(\widehat h\) is an \(\widehat N\)-residual.  Therefore
\[
 \widehat B^{\mathsf T}\widehat k=0. \tag{117}
\]
Using \(d=Bw\), conditioning on \((Y,X)\), subtracting (117), and using \(gw_j=\pi_jf_j\) gives
\[
 E_P(\widehat k\mid X=x)
 =\sum_{j=1}^K\pi_j(x)\int_{\mathcal Y}f_j(y\mid x)
 \{B_{\cdot j}(x)-\widehat B_{\cdot j}(x)\}^{\mathsf T}
 \widehat k(y,x)d\nu(y). \tag{118}
\]
Let \(\Delta_A=M_A-\bar M_A\) and \(\Delta_Z=M_Z-\bar M_Z\).  Columnwise tensor multiplication yields
\[
 B_{\cdot j}-\widehat B_{\cdot j}
 =\Delta_{A,\cdot j}\otimes M_{Z,\cdot j}
 +M_{A,\cdot j}\otimes\Delta_{Z,\cdot j}
 -\Delta_{A,\cdot j}\otimes\Delta_{Z,\cdot j}. \tag{119}
\]
Theorems \(\mathrm{thm{:}2\mbox{-}tangent\mbox{-}completeness}\) and
\(\mathrm{thm{:}3\mbox{-}finite\mbox{-}projector}\) imply
\[
 E_P(k\mid A,J,X)=0,\qquad E_P(k\mid Z,J,X)=0. \tag{120}
\]
For example, if \(M_A(x)_{aj}>0\), three-view independence and (120) give
\[
 \int_{\mathcal Y}f_j(y\mid x)
 \sum_zM_Z(x)_{zj}k_{az}(y,x)d\nu(y)=0. \tag{121}
\]
If \(M_A(x)_{aj}=0\), exact recovery on \(\mathcal E_{\mathrm{clean},n}\) gives
\(\Delta_A(x)_{aj}=0\).  Thus the term in (119) linear in \(\Delta_A\) vanishes after replacing \(\widehat k\) by \(k\), cell by cell.  The symmetric assertion holds for the term linear in \(\Delta_Z\).

Write \(\widehat k=k+(\widehat k-k)\) in the two linear terms and retain \(\widehat k\) in the product term.  The true and fitted positivity and singular-value margins, bounded outcome, and
\(\|\widehat G_{\tau_G}^{\dagger}\|_{\mathrm{op}}\leq2/\varepsilon\)
bound \(k\) and \(\widehat k\) uniformly.  Conditional Cauchy--Schwarz, (114), and integration over \(P_X\) yield
\[
 \|E_P\widehat k\|
 \leq C_{\mathrm{reg}}
 \left[(r_A^c+r_Z^c)\|\widehat k-k\|_2+r_A^cr_Z^c\right]. \tag{122}
\]
Equations (113)--(116) and (122) prove every assertion.
\end{proof}
\par\smallskip\noindent \textit{Justification.} Thresholding and renormalization recover true zeros exactly on the screening event, making the two linear measurement terms cancel without beta-min. \textit{Gap.} A diagnostic-only leakage allowance does not yield the product-only remainder needed for flexible root-\(n\) inference. \textit{Consumer.} It removes additive unknown-support leakage from the projected-score bias.

\begin{lemma}[lem:cleaned-support-rate-window]\label{lem:cleaned-support-rate-window}
Put \(D_0=d_X+1\), \(\rho_n=n^{-\beta/(2\beta+D_0)}\), and \(\zeta_n=\rho_n\sqrt{\log n}\).  If \(\beta>D_0/2\) and a label-aligned raw measurement fit satisfies
\[
 r_{A,\infty}+r_{Z,\infty}=O_P(\zeta_n),
 \qquad r_A+r_Z=O_P(\rho_n),
\]
then the deterministic choice
\[
 \tau_{S,n}=\{\zeta_n n^{-1/4}\}^{1/2}
\]
satisfies
\[
 \zeta_n=o(\tau_{S,n}),
 \qquad \tau_{S,n}=o(n^{-1/4}),
 \qquad \kappa_{\mathrm{supp},n}=o_P(1).
\]
For the threshold-and-renormalize cleaning, it also gives
\[
 r_A^c+r_Z^c=O_P(\rho_n+\tau_{S,n})=o_P(n^{-1/4}),
\]
and exact recovery of every true-zero measurement cell with probability tending to one, without a beta-min condition on positive cells.
\end{lemma}
\begin{proof}
Let \(\gamma=\beta/(2\beta+D_0)\).  The condition \(2\beta>D_0\) is equivalent to \(\gamma>1/4\).  Therefore
\[
 n^{1/4}\zeta_n=n^{1/4-\gamma}\sqrt{\log n}\longrightarrow0.  \tag{87}
\]
For the displayed choice of \(\tau_{S,n}\), direct division gives
\[
 \frac{\zeta_n}{\tau_{S,n}}
 =\frac{\tau_{S,n}}{n^{-1/4}}
 =\{n^{1/4}\zeta_n\}^{1/2}\longrightarrow0.                   \tag{88}
\]
This proves the two deterministic rate relations.  Since
\((r_{A,\infty}+r_{Z,\infty})/\zeta_n=O_P(1)\), multiplication by the first ratio in (88) yields
\[
 \kappa_{\mathrm{supp},n}
 =\frac{r_{A,\infty}+r_{Z,\infty}}{\tau_{S,n}}=o_P(1).         \tag{89}
\]
By definition, \(r_A^c+r_Z^c=r_A+r_Z+2\tau_{S,n}\), so the assumed \(L^2\) rate and (88) imply
\[
 r_A^c+r_Z^c=O_P(\rho_n+\tau_{S,n})=o_P(n^{-1/4}),              \tag{90}
\]
where \(\rho_n=o(n^{-1/4})\) follows from \(\gamma>1/4\).  Finally, (89) and the already proved support-cleaning lemma give exact true-zero recovery with probability tending to one.  That lemma's deterministic deletion bound controls any deleted positive cell by \(C_{\mathrm{reg}}\tau_{S,n}\), so no positive-cell lower bound has been used.
\end{proof}
\par\smallskip\noindent \textit{Justification.} It proves that the screening and root-\(n\) cleaning requirements have a nonempty explicit threshold window under the proposed uniform rate. \textit{Gap.} The earlier diagnostic formulation did not establish compatibility of exact support recovery with a second-order threshold cost. \textit{Consumer.} It separates the closed support algebra from the genuinely open spline-likelihood rate proof.

\section{Interpretation}

Canonicality is a property of the full observed-data DQM tangent space on the declared two-sided pathwise-interior, quantitatively separated model.  Implementation uses a different but compatible argument: low raw fitted measurement probabilities are deleted and each column is renormalized.  With high probability this deletes every true-zero cell, while deletion of arbitrarily small positive cells costs only \(O(\tau_{S,n})\); no positive-cell beta-min premise is used.  Consequently the two measurement-linear bias terms cancel against the true tangent constraints and only products of cleaned errors remain.  The raw hatted matrices are screening inputs; all bridges and projection objects use the barred matrices.

\section*{Technical internal limitation (diagnostic only)}

The exact differential inverse and canonical gradient are proved only on the declared fixed-\(K\), common-support, two-sided pathwise-interior stratum with quantitative measurement, posterior, and projection-Gram gaps.  This is a separated specialization, not the full class studied by Zhou and Tchetgen Tchetgen or by Bian, Zhou, and Cui.  Support cleaning additionally requires \(r_{A,\infty}+r_{Z,\infty}=o_P(\tau_{S,n})\) and, for root-\(n\) inference with a training-fold-random threshold, \(\tau_{S,n}=o_P(n^{-1/4})\).  The frozen H"older conditions alone do not prove the local curvature, spline boundary, localization, or uniform-rate properties of the specifically prescribed global product-likelihood estimator; that construction remains open.

\section{Honest scope}

Exact identification, tangent completeness, the finite projector, canonicality, strict variance reduction, the binary comparison, and the continuous witness remain claims about the declared quantitatively separated fixed-\(K\) model.  Efficient cross-fitted inference is conditional on a coherent support-cleaned fit satisfying \(\kappa_{\mathrm{supp},n}=o_P(1)\), \(\tau_{S,n}=o_P(n^{-1/4})\), \(R_\Sigma=o_P(1)\), and \(\Delta_n=o_P(n^{-1/2})\).  The algorithm uses neither unknown true measurement supports nor an oracle Gram rank, and it imposes no beta-min condition on positive measurement cells.  A complete rate proof for the named spline product likelihood is not claimed.  The maternal-smoking application is prospective because the published six-bin design does not itself verify paired-proxy conditional independence or the quantitative separation margins.

\begin{thebibliography}{99}

  \bibitem{ZhouTchetgenTchetgen2026} Ying Zhou and Eric J. Tchetgen Tchetgen (2026). Causal Inference with a Hidden Treatment. arXiv:2405.09080, version 4.

  \bibitem{BianZhouCui2026} Zeyu Bian, Ying Zhou, and Yifan Cui (2026). Learning from the Unseen: Offline Reinforcement Learning with Hidden Actions. arXiv:2607.25241, version 1.

  \bibitem{GuoOgburnShpitser2026} Helen Guo, Elizabeth L. Ogburn, and Ilya Shpitser (2026). Comparing Two Proxy Methods for Causal Identification. arXiv:2512.00175, version 3.

  \bibitem{AllmanMatiasRhodes2009} Elizabeth S. Allman, Catherine Matias, and John A. Rhodes (2009). Identifiability of Parameters in Latent Structure Models with Many Observed Variables. Annals of Statistics, 37(6A), 3099--3132.

  \bibitem{HuSchennach2008} Yingyao Hu and Susanne M. Schennach (2008). Instrumental Variable Treatment of Nonclassical Measurement Error Models. Econometrica, 76(1), 195--216.

  \bibitem{Deaner2023} Ben Deaner (2023). Controlling for Latent Confounding with Triple Proxies. arXiv:2204.13815, revised manuscript.

  \bibitem{TsiatisMa2004} Anastasios A. Tsiatis and Yanyuan Ma (2004). Locally Efficient Semiparametric Estimators for Functional Measurement Error Models. Biometrika, 91(4), 835--848.

  \bibitem{Tsybakov2009} Alexandre B. Tsybakov (2009). Introduction to Nonparametric Estimation. Springer.

  \bibitem{BonhommeJochmansRobin2016} St\'ephane Bonhomme, Koen Jochmans, and Jean-Marc Robin (2016). Non-parametric Estimation of Finite Mixtures from Repeated Measurements. Journal of the Royal Statistical Society, Series B, 78(1), 211--229.

  \bibitem{BennettEtAl2023} Andrew Bennett, Nathan Kallus, Xiaojie Mao, Whitney Newey, Vasilis Syrgkanis, and Masatoshi Uehara (2023). Inference on Strongly Identified Functionals of Weakly Identified Functions. Proceedings of Machine Learning Research, 195.

  \bibitem{GhassamiEtAl2026} AmirEmad Ghassami, Yunshu Zhang, Ilya Shpitser, and Eric J. Tchetgen Tchetgen (2026). Partial Identification of Causal Effects Using Proxy Variables. arXiv:2304.04374, version 4.

  \bibitem{CobzaruEtAl2024} Raluca Cobzaru, Roy Welsch, Stan Finkelstein, Kenney Ng, and Zach Shahn (2024). Bias Formulas for Violations of Proximal Identification Assumptions in a Linear Structural Equation Model. Journal of Causal Inference, 12(1), 1--34.

  \bibitem{BodoryBusshoffLechner2022} Philipp Bodory, Jan Busshoff, and Michael Lechner (2022). High Resolution Treatment Effects Estimation: Uncovering Effect Heterogeneities with the Modified Causal Forest. Entropy, 24(8), 1039.

\end{thebibliography}

\end{document}